% --- Section: 00_meta.tex ---
\title{\textbf{Adaptive SVD Attention: Energy-Thresholded Low-Rank Factorization with Residual Adapters for Multi-Architecture Transformers}}

\author{Anonymous authors}


% --- Section: 01_intro.tex ---

\section{Introduction}

\subsection{Motivation}

The four attention projections $W_Q, W_K, W_V, W_O$ in a standard transformer 
layer account for the majority of attention parameters. For DistilBERT 
($d = 768$), each projection contains $768 \times 768 = 589{,}824$ parameters; 
the four together contain ~2.36M parameters per layer.

\textbf{Low-rank factorization} replaces $W \in \mathbb{R}^{d_\text{out} \times d_\text{in}}$ 
with the product $U V$ where $U \in \mathbb{R}^{d_\text{out} \times r}$ and 
$V \in \mathbb{R}^{r \times d_\text{in}}$. Parameter count drops from 
$d_\text{in} d_\text{out}$ to $r (d_\text{in} + d_\text{out})$. For 
$768 \times 768$ at $r = 128$, this is a \textbf{3$\times$ reduction} in attention 
parameters.

The \textbf{rank choice} is the central design decision. Truncated SVD gives the 
best rank-$r$ approximation in Frobenius norm (Eckart--Young), but the theorem 
says nothing about what $r$ should be. Existing approaches fall into three 
categories:

\begin{enumerate}
\item \textbf{Fixed global rank.} Apply the same $r$ to every projection. Simple, but
  ignores that projections have very different \textbf{spectral decay rates}.
  \textbf{Figure 1.1} shows singular value spectra of two DistilBERT projections:
  one decays rapidly (~20 effective rank), another decays slowly (~200).

\item \textbf{Validation-driven escalation.} Start low, increase rank if Frobenius
  error exceeds a threshold. Requires multiple SVDs per layer.

\item \textbf{Activation-aware rank.} Weight the SVD input space by empirical
  activation statistics (Fisher diagonal, covariance). Requires calibration
  data.

\end{enumerate}
We propose a \textbf{fourth approach} that is self-contained: use the singular 
values themselves to determine rank. Specifically, threshold **cumulative 
squared energy**:

\[
\mathcal{E}(r) = \frac{\sum\_{i=1}^{r} \sigma\_i^2}{\sum\_{j=1}^{\min(d\_\text{in}, d\_\text{out})} \sigma\_j^2}, \quad
r^*(\tau) = \min\{r : \mathcal{E}(r) \geq \tau\}
\tag{1.1}
\]

with default $\tau = 0.92$. Fast-decaying projections get 
small $r$, slow-decaying ones get large $r$, all determined automatically.

\subsection{Problem Statement}

Given a pretrained attention module with projections 
$W_P \in \mathbb{R}^{d_\text{out}^{(P)} \times d_\text{in}^{(P)}}$ for 
$P \in \{Q, K, V, O\}$, and an energy threshold $\tau \in (0, 1)$, we seek 
per-projection factorizations $W_P \approx U_P V_P$ such that:

\begin{enumerate}
\item \textbf{Energy preservation:} $\mathcal{E}_P(r_P) \geq \tau$ for each $P$.
\item \textbf{Parameter reduction:} $r_P (d_\text{in}^{(P)} + d_\text{out}^{(P)}) < d_\text{in}^{(P)} d_\text{out}^{(P)}$.
\item \textbf{No base retraining:} $U_P, V_P$ are frozen; only adapters train.
\item \textbf{No calibration data:} rank selected purely from the weight's spectrum.
\item \textbf{Multi-architecture support:} flat, nested, causal, T5, GQA, RoPE.
\item \textbf{Initial forward fidelity:} with adapters zero-initialized, the initial
  forward pass matches the truncation approximation exactly.

\end{enumerate}
\subsection{Our Approach}

\textbf{Adaptive SVD Attention} combines four design choices:

\begin{enumerate}
\item \textbf{Cumulative energy rank.} For each projection, compute SVD, then select
  the smallest rank $r$ where cumulative squared energy $\mathcal{E}(r) \geq \tau$.

\item \textbf{Symmetric factor distribution.} Split $\Sigma^{1/2}$ between factors:

\end{enumerate}
\[
U\_r = U\_{:, :r} \Sigma\_{:r, :r}^{1/2}, \quad
V\_r = \Sigma\_{:r, :r}^{1/2} V\_{:r, :}^\top
\tag{1.2}
\]

This keeps $\|U_r\|_2 = \|V_r\|_2 = \sqrt{\sigma_\text{max}}$, avoiding FP16 overflow.

\begin{enumerate}
\item \textbf{Residual adapters.} Add zero-initialized $B A$ with $B = 0$ at init. The
  adapter contributes zero to the initial forward, giving a stable starting
  point for recovery.

\item \textbf{Universal architecture support.} The injection handles flat (DistilBERT),
  nested (RoBERTa/BERT), fused Conv1D (GPT-2), GQA (Qwen), RoPE (LLaMA),
  and T5 relative attention bias.

\end{enumerate}
\subsection{Contributions}

\begin{itemize}
\item \textbf{Energy-thresholded rank selection.} A self-contained, calibration-free
  criterion for per-layer rank based on cumulative singular value energy.
  Section 4.2 formalizes this.
\item \textbf{Symmetric SVD + residual adapter.} A factorization that keeps both
  factors balanced in magnitude and adds zero-initialized residuals for
  recovery. Section 4.3--4.4.
\item \textbf{Broad architecture support.} The method works on flat, nested, and causal
  attention, with optional T5 bias, GQA, and RoPE. Section 5.
\item \textbf{Systematic evaluation.} We benchmark on 3 models on SST-2 with ablations
  over $\tau$, residual rank, and recovery steps. Section 7.
\item \textbf{Honest reporting.} We report which models benefit most, which have
  limitations, and what we tried that failed (Appendix D).
\item \textbf{Full reproducibility.} Code, checkpoints, and evaluation scripts released
  under MIT.

\end{itemize}
\subsection{Paper Organization}

Section 2 reviews related work in low-rank compression, activation-aware 
methods, and parameter-efficient fine-tuning. Section 3 provides background 
on SVD, energy thresholds, and LoRA. Section 4 presents the method. Section 5 
describes integration with real transformer architectures. Section 6 details 
experimental setup. Section 7 reports results. Section 8 analyzes the 
per-layer rank distribution and ablation contributions. Section 9 discusses 
practical usage. Section 10 lists limitations. Section 11 concludes.



\subsection{Notation}


% --- Section: 02_related.tex ---

\section{Related Work}

\subsection{Low-Rank Matrix Factorization}

\textbf{Classical SVD.} The Eckart--Young theorem (Eckart \& Young, 1936) states 
that the best rank-$r$ approximation of $W$ in Frobenius norm is the truncated 
SVD $W_r = U_{:, :r} \Sigma_{:r, :r} V_{:r, :}^\top$. Adaptive SVD builds 
on this but selects $r$ via an energy threshold rather than a validation 
loop or fixed value.

\textbf{Rank selection heuristics.} Classical rank selection includes AIC 
(Akaike, 1974), BIC (Schwarz, 1978), and cross-validation. For SVD-based 
compression specifically, prior work uses reconstruction error thresholds 
(A3), fixed budget ratios (SDRM), or optimal hard thresholding on 
singular values under noise (Gavish \& Donoho, 2014). Our criterion is a 
\textbf{cumulative energy threshold}, distinct from per-singular-value thresholds.

\textbf{Randomized SVD.} Halko et al. (2011) provide a randomized algorithm with 
$\mathcal{O}(d r \log k)$ complexity. For our layer sizes ($\le$3072	imes$768), exact 
SVD is faster and simpler.

\subsection{Activation-Aware Compression}

\textbf{AWQ.} Lin et al. (2023) observe that activation magnitude correlates with 
weight importance and use it to scale quantization. Applied to low-rank 
factorization, this becomes the A3 method (Fisher diagonal weighting).

\textbf{SmoothQuant.} Xiao et al. (2023) migrate activation scale to weights for 
quantization. Analogous to AASVD's covariance weighting for low-rank.

\textbf{A3.} Uses Fisher diagonal $F\_i = \mathbb{E}[x\_i^2]$ to weight the SVD input 
space. Requires calibration. Adaptive SVD is calibration-free.

\textbf{AASVD.} Uses full covariance $C\_X = \mathbb{E}[x x^\top]$. Requires 
streaming EMA accumulation. Adaptive SVD is calibration-free.

\textbf{Key distinction.} Adaptive SVD does \textbf{not} use activation statistics. Its 
rank criterion is intrinsic to the weight matrix.

\subsection{Parameter-Efficient Fine-Tuning}

\textbf{LoRA.} Hu et al. (2021) add trainable rank-$r\_L$ updates $B A$ with $B = 0$. 
Adaptive SVD uses the same residual structure, integrated with a frozen SVD base.

\textbf{AdaLoRA.} Zhang et al. (2023) adaptively allocate LoRA rank per layer based 
on singular value importance. Our adaptation is on the \textbf{base factorization}, 
not the adapter.

\textbf{PiSSA.} Meng et al. (2024) initialize LoRA adapters with principal singular 
components. Adaptive SVD uses the top-$r$ components as the \textbf{frozen base} 
and zero-initialized LoRA as the residual. These are complementary strategies.

\textbf{OFT.} Qiu et al. (2023) use orthogonal fine-tuning. Different primitive.

\subsection{Attention-Specific Compression}

\textbf{Head pruning.} Voita et al. (2019) prune entire heads. Adaptive SVD preserves 
head count and compresses per-head projections via low-rank factorization.

\textbf{GQA.} Ainslie et al. (2023) reduce KV head count during pretraining. Adaptive 
SVD preserves head count but compresses per-head projection matrices.

\textbf{KV cache compression.} H2O (Zhang et al., 2023), StreamingLLM (Xiao et al., 
2024) reduce cache size at inference. Different target.

\subsection{Position of Our Work}

\begin{table}[H]
\centering
\small
\begin{tabular}{llllll}
\toprule
\textbf{Method} & \textbf{Calib} & \textbf{Weighted} & \textbf{Rank} & \textbf{Per-layer rank} & \textbf{T5 bias} \\
\midrule
SVD fixed rank \& No \& No \& Fixed \& No \& No \\
SDRM \& No \& No \& Budget \& No \& No \\
A3 \& Yes \& Fisher diag \& Escalation \& Partial \& No \\
AASVD \& Yes \& Full cov \& Fixed \& No \& No \\
\textbf{Adaptive SVD} & \textbf{No} & \textbf{No} & \textbf{Energy} & \textbf{Yes} & \textbf{Yes} \\
\bottomrule
\end{tabular}
\end{table}

\textbf{Key distinction.} Adaptive SVD is the only method that:
\begin{enumerate}
\item Requires no calibration data.
\item Selects rank per-projection via an energy criterion.
\item Preserves T5 relative attention bias.

\end{enumerate}
\subsection{Comparison Table}


% --- Section: 03_background.tex ---

\section{Background}

\subsection{Singular Value Decomposition}

For any matrix $W \in \mathbb{R}^{m \times n}$, the SVD is:

\[
W = U \Sigma V^\top
\tag{3.1}
\]

where $U \in \mathbb{R}^{m \times m}$, $V \in \mathbb{R}^{n \times n}$ are 
orthogonal, and $\Sigma = \mathrm{diag}(\sigma_1, \dots, \sigma_{\min(m,n)})$ 
with $\sigma_1 \geq \sigma_2 \geq \dots \geq 0$.

\textbf{Truncated SVD.} The best rank-$r$ approximation in Frobenius norm is:

\[
W\_r = U\_{:, :r} \Sigma\_{:r, :r} V\_{:r, :}^\top
\tag{3.2}
\]

with $\|W - W_r\|_F^2 = \sum_{i > r} \sigma_i^2$.

\textbf{Symmetric factorization.} Distributing $\Sigma^{1/2}$ between factors:

\[
U\_r = U\_{:, :r} \Sigma\_{:r, :r}^{1/2}, \quad
V\_r = \Sigma\_{:r, :r}^{1/2} V\_{:r, :}^\top
\tag{3.3}
\]

ensures $\|U_r\|_2 = \|V_r\|_2 = \sigma_1^{1/2}$, giving balanced magnitudes 
and improved numerical stability for FP16 computation.

\subsection{Cumulative Squared Energy}

Define the cumulative squared energy at rank $r$:

\[
\mathcal{E}(r) = \frac{\sum\_{i=1}^{r} \sigma\_i^2}{\sum\_{j=1}^{\min(m,n)} \sigma\_j^2}
\tag{3.4}
\]

Properties:
\begin{itemize}
\item $\mathcal{E}(0) = 0$
\item $\mathcal{E}$ is monotone increasing
\item $\mathcal{E}(\min(m, n)) = 1$
\item $\mathcal{E}(r)$ measures the fraction of $\|W\|_F^2$ captured by the top-$r$ directions

\end{itemize}
\textbf{Energy-thresholded rank.} Given $\tau \in (0, 1)$:

\[
r^*(\tau) = \min\{r : \mathcal{E}(r) \geq \tau\}
\tag{3.5}
\]

At this rank:

\[
\|W - W\_{r^*}\|\_F^2 = (1 - \mathcal{E}(r^*)) \|W\|\_F^2 \leq (1 - \tau) \|W\|\_F^2
\tag{3.6}
\]

so relative Frobenius error is $\leq \sqrt{1 - \tau}$.

\subsection{Spectral Decay in Attention Weights}

Attention projection matrices exhibit heterogeneous spectral decay:

\begin{itemize}
\item \textbf{Query projections:} Moderate decay. Effective rank ~40--80\% of full.
\item \textbf{Key projections:} Slower decay (keys spread information across dimensions).
\item \textbf{Value projections:} Faster decay (values have lower effective rank).
\item \textbf{Output projections:} Moderate decay; depends on head-combining structure.

\end{itemize}
\textbf{Figure 3.1} shows singular value spectra for DistilBERT layer 0 vs layer 3. 
The difference in decay rate justifies per-layer rank selection.

\subsection{LoRA-Style Residual Adapters}

LoRA (Hu et al., 2021) parameterizes a weight update as:

\[
\Delta W = \frac{\alpha}{r\_L} B A
\tag{3.7}
\]

with $A \in \mathbb{R}^{r_L \times d_\text{in}}$, $B \in \mathbb{R}^{d_\text{out} \times r_L}$. 
The initialization $B = 0$ ensures $\Delta W = 0$ at step 0.

\textbf{Applied to factorized layers.} The full forward:

\[
y = x (U V)^\top + \frac{\alpha}{r\_L} x (B A)^\top
\tag{3.8}
\]

At initialization, $y = x (U V)^\top$, matching the frozen base.

\subsection{T5 Relative Attention Bias}

T5 (Raffel et al., 2020) adds a learnable bias to attention scores based on 
the bucketed relative position of query and key:

\[
b(i, j) = \text{Embed}[\text{bucket}(j - i)]
\tag{3.9}
\]

The bucket function is:

\[
\text{bucket}(d) = \begin{cases}
d & \text{if } |d| \textless{} \text{max\_exact} \\
\text{max\_exact} + \text{sign}(d) \cdot \lfloor \log(|d| / \text{max\_exact}) \cdot \text{scale} \rfloor & \text{otherwise}
\end{cases}
\tag{3.10}
\]

with $\text{max\_exact} = 8$, $\text{num\_buckets} = 32$. The bias is added 
to attention scores before softmax. Adaptive SVD preserves this bias.

\subsection{Grouped-Query Attention (GQA)}

Modern models (Qwen, LLaMA) use GQA where $H_{kv} < H$:

\[
K, V \in \mathbb{R}^{B \times L \times H\_{kv} d\_h}, \quad H\_\text{repeat} = H / H\_{kv}
\tag{3.11}
\]

The KV heads are repeated via \texttt{repeat\_kv} before attention. Adaptive SVD 
preserves GQA structure.

\subsection{Rotary Position Embeddings (RoPE)}

RoPE (Su et al., 2021) rotates $Q, K$ by position-dependent angles:

\[
\tilde{Q}\_i = Q\_i \cos(\theta\_i) + \mathrm{rot}(Q\_i) \sin(\theta\_i)
\tag{3.12}
\]

Adaptive SVD applies RoPE after the factorized $Q, K$ projections.

\subsection{Notation Summary}


% --- Section: 04_method.tex ---

\section{Method}

\subsection{Overview}

Adaptive SVD compresses each attention projection in four steps:

\begin{enumerate}
\item \textbf{Full SVD} ( 4.2): compute $U \Sigma V^\top$ in FP32.
\item \textbf{Energy-thresholded rank} ( 4.3): find smallest $r$ with $\mathcal{E}(r) \geq \tau$.
\item \textbf{Symmetric factorization} ( 4.4): set $U_r = U_{:, :r} \Sigma^{1/2}$,
  $V_r = \Sigma^{1/2} V_{:r, :}^\top$.
\item \textbf{Residual adapter} ( 4.5): attach zero-initialized $B A$ for recovery.

\end{enumerate}
\textbf{Figure 4.1} shows the pipeline. There is no calibration phase.

\begin{lstlisting}
    +------------------------------------------+
    |  Pretrained attention module A           |
    |  with projections W\_Q, W\_K, W\_V, W\_O     |
    +------------------+-----------------------+
                       |
                        
    +------------------------------------------+
    |  For each projection P:                  |
    |    Phase 1: SVD(W\_P) $\to$ U\_P  \_P V\_P^T     |
    |    Phase 2: Find r\_P via $\tau$               |
    |    Phase 3: Truncate + distribute  ^{1/2}|
    |    Phase 4: Add zero-init BA adapter     |
    +------------------+-----------------------+
                       |
                        
    +------------------------------------------+
    |  Compressed attention A' with:           |
    |    - Frozen U\_P, V\_P for each P          |
    |    - Trainable A\_P, B\_P for each P       |
    +------------------------------------------+
\end{lstlisting}

\subsection{Phase 1: Full SVD}

For each projection weight $W_P \in \mathbb{R}^{d_\text{out}^{(P)} \times d_\text{in}^{(P)}}$:

\[
W\_P = U\_P \Sigma\_P V\_P^\top
\tag{4.1}
\]

We compute this in \textbf{FP32} regardless of the model dtype. FP16 SVD can 
produce NaN for ill-conditioned matrices.

\textbf{Complexity.} $\mathcal{O}(d_\text{in} d_\text{out} \min(d_\text{in}, d_\text{out}))$. 
For $768 \times 768$, this is ~450M FLOPs, ~50 ms on CPU, ~5 ms on GPU.

\subsection{Phase 2: Energy-Thresholded Rank}

\subsubsection{Definition}

Given singular values $\sigma_1 \geq \sigma_2 \geq \dots$, define:

\[
s\_i = \sigma\_i^2, \quad \mathcal{E}(r) = \frac{\sum\_{i=1}^{r} s\_i}{\sum\_{j=1}^{\min(d\_\text{in}, d\_\text{out})} s\_j}
\tag{4.2}
\]

The rank is:

\[
r\_P = \min\{r : \mathcal{E}(r) \geq \tau\}
\tag{4.3}
\]

\subsubsection{Implementation}

\begin{lstlisting}
def \_compute\_energy\_rank(self, S, max\_rank, min\_rank=16):
    s\_squared = S ** 2
    total\_energy = s\_squared.sum()
    if total\_energy \textless{}= 1e-12:
        return max(1, min(min\_rank, max\_rank))
    
    cum\_energy = torch.cumsum(s\_squared, dim=0) / total\_energy
    idx = torch.where(cum\_energy \textgreater{}= self.energy\_threshold)[0]
    r = int(idx[0].item()) + 1 if len(idx) \textgreater{} 0 else max\_rank
    return max(1, min(max(r, min\_rank), max\_rank))
\end{lstlisting}

\textbf{Complexity.} $\mathcal{O}(\min(d_\text{in}, d_\text{out}))$ via \texttt{cumsum}.

\subsubsection{Minimum Rank Constraint}

We impose $r_P \geq r_\text{min}$ (default 16) to prevent numerical 
instability in very low-rank layers. This bounds $r_P$ from below, but for 
$\tau = 0.92$ and typical attention weights, the natural rank is far above 
$r_\text{min}$.

\subsubsection{Maximum Rank Cap}

We cap $r_P \leq \min(d_\text{in}, d_\text{out})$. In practice, $\tau < 1.0$ 
ensures $r_P < \min(d_\text{in}, d_\text{out})$, but the cap protects against 
edge cases.

\subsection{Phase 3: Symmetric Factorization}

\subsubsection{Factor Construction}

Given $r_P$, set:

\[
U\_{P, r} = U\_{P, :, :r\_P} \Sigma\_{P, :r\_P, :r\_P}^{1/2}, \quad
V\_{P, r} = \Sigma\_{P, :r\_P, :r\_P}^{1/2} V\_{P, :r\_P, :}^\top
\tag{4.4}
\]

Then $U_{P, r} V_{P, r} = W_{P, r_P}$ (the truncated SVD).

\subsubsection{Why Symmetric?}

Three reasons:

\begin{enumerate}
\item \textbf{Numerical stability.} Both factors have magnitude $\sqrt{\sigma_\text{max}}$,
  avoiding FP16 overflow.

\item \textbf{Balanced gradient flow.} During recovery, adapters interact symmetrically
  with both factors.

\item \textbf{Orthogonality.} When singular values are distinct, columns of $U_{P, r}$
  and rows of $V_{P, r}$ are orthogonal.

\end{enumerate}
\subsubsection{The Factorized Forward Pass}

The compressed projection computes:

\[
\text{Proj}\_P(x) = (x V\_{P, r}^\top) U\_{P, r}^\top + b\_P
\tag{4.5}
\]

\textbf{Parameter count:}

\begin{table}[H]
\centering
\small
\begin{tabular}{ll}
\toprule
\textbf{Configuration} & \textbf{Parameters} \\
\midrule
Original & $d_\text{in}^{(P)} d_\text{out}^{(P)}$ \\
Adaptive SVD & $r_P (d_\text{in}^{(P)} + d_\text{out}^{(P)})$ \\
\bottomrule
\end{tabular}
\end{table}

\textbf{Reduction factor:} $\frac{d_\text{in} d_\text{out}}{r_P (d_\text{in} + d_\text{out})}$.

\subsection{Phase 4: Residual Adapter}

\subsubsection{The Residual}

The truncation residual is:

\[
\Delta W\_P = W\_P - U\_{P, r} V\_{P, r}
\tag{4.6}
\]

Its Frobenius norm is $\sqrt{(1 - \mathcal{E}(r_P))} \cdot \|W_P\|_F \leq \sqrt{1 - \tau} \|W_P\|_F$. 
At $\tau = 0.92$, this is ~28\% of $\|W_P\|_F$.

\subsubsection{Adapter Parameterization}

We approximate the residual with a rank-$r_L$ factor:

\[
\Delta W\_P \approx B\_P A\_P, \quad
A\_P \in \mathbb{R}^{r\_L \times d\_\text{in}}, \quad
B\_P \in \mathbb{R}^{d\_\text{out} \times r\_L}
\tag{4.7}
\]

\subsubsection{Initialization}

\[
B\_P = 0, \quad A\_P \sim \mathcal{N}(0, 1/\sqrt{r\_L})
\tag{4.8}
\]

The zero initialization of $B_P$ ensures $\Delta W_P = B_P A_P = 0$ at step 0.

\subsubsection{Full Forward Pass}

\[
\text{Proj}\_P(x) = (x V\_{P, r}^\top) U\_{P, r}^\top + (x A\_P^\top) B\_P^\top + b\_P
\tag{4.9}
\]

\textbf{Parameter overhead:} $r_L (d_\text{in} + d_\text{out})$ per projection.

\subsection{The Full Algorithm}

\textbf{Algorithm 1:} Adaptive SVD compression for one attention module
\begin{lstlisting}
Input:  pretrained attention A with projections P $\in$ {Q, K, V, O}
        energy threshold $\tau$ $\in$ (0, 1)
        residual rank r\_L
        minimum rank r\_min
Output: compressed attention A'

1:  for P in {Q, K, V, O} do
2:      W\_P $\leftarrow$ weight(A.P)                              Extract
3:      U\_P,  \_P, V\_P^T $\leftarrow$ svd(W\_P)                     Full SVD in FP32
4:      s $\leftarrow$  \_P^2
5:      E $\leftarrow$ cumsum(s) / sum(s)                        Cumulative energy
6:      idx $\leftarrow$ first index where E $\ge$ $\tau$
7:      r\_P $\leftarrow$ max(idx + 1, r\_min)                     Enforce min rank
8:
9:       \_r^{1/2} $\leftarrow$ sqrt( \_P[:r\_P])
10:     U\_r $\leftarrow$ U\_P[:, :r\_P] @ diag( \_r^{1/2})          Symmetric split
11:     V\_r $\leftarrow$ diag( \_r^{1/2}) @ V\_P^T[:r\_P, :]
12:
13:     B\_P $\leftarrow$ zeros(d\_out, r\_L)                       Zero init
14:     A\_P $\leftarrow$ kaiming\_uniform(r\_L, d\_in)
15:
16:     replace A.P with AdaptiveSVDLinear(U\_r, V\_r, A\_P, B\_P)
17: end for
18:
19: freeze(U\_r, V\_r) for all P                        Base frozen
20: unfreeze(A\_P, B\_P) for all P                      Adapters trainable
21:
22: return A'
\end{lstlisting}

\subsection{Complexity Analysis}

\subsubsection{Time Complexity}

\begin{table}[H]
\centering
\small
\begin{tabular}{ll}
\toprule
\textbf{Phase} & \textbf{Time} \\
\midrule
SVD per projection & $\mathcal{O}(d_\text{in} d_\text{out} \min(d_\text{in}, d_\text{out}))$ \\
Rank selection & $\mathcal{O}(\min(d_\text{in}, d_\text{out}))$ \\
Factor truncation & $\mathcal{O}(r_P (d_\text{in} + d_\text{out}))$ \\
Adapter init & $\mathcal{O}(r_L (d_\text{in} + d_\text{out}))$ \\
\textbf{Per projection} & $\mathcal{O}(d_\text{in} d_\text{out} \min(d_\text{in}, d_\text{out}))$ \\
\textbf{Per attention module} & $\mathcal{O}(4 \cdot d^3)$ for square $d \times d$ \\
\textbf{Total ($N_l$ layers)} & $\mathcal{O}(4 N_l d^3)$ \\
\bottomrule
\end{tabular}
\end{table}

For DistilBERT ($N_l = 6$, $d = 768$): ~30 seconds total on GPU.

\subsubsection{Space Complexity}

Peak memory per projection: $O(d^2)$ for the SVD. All covariances 
and calibration states are avoided.

\subsubsection{Compression Time Comparison}

\begin{table}[H]
\centering
\small
\begin{tabular}{llll}
\toprule
\textbf{Method} & \textbf{Calibration} & \textbf{SVD} & \textbf{Total} \\
\midrule
A3 & ~10s & ~20s & ~30s \\
AASVD & ~15s & ~20s & ~35s \\
\textbf{Adaptive SVD} & \textbf{0s} & \textbf{~15s} & \textbf{~15s} \\
\bottomrule
\end{tabular}
\end{table}

\subsection{Parameter Analysis}

For a square $d \times d$ projection:

\begin{table}[H]
\centering
\small
\begin{tabular}{ll}
\toprule
\textbf{Configuration} & \textbf{Parameters} \\
\midrule
Original & $d^2$ \\
Adaptive SVD & $2 r_P d$ \\
Adaptive SVD + LoRA($r_L$) & $2 d (r_P + r_L)$ \\
\bottomrule
\end{tabular}
\end{table}

\textbf{Reduction factor (with adapter):}

\[
\frac{d^2}{2 d (r\_P + r\_L)} = \frac{d}{2(r\_P + r\_L)}
\tag{4.10}
\]

\textbf{Example.} $d = 768$, $r_P = 128$, $r_L = 16$:

\begin{table}[H]
\centering
\small
\begin{tabular}{lll}
\toprule
\textbf{Configuration} & \textbf{Parameters} & \textbf{Ratio} \\
\midrule
Original & 589,824 & 1.00$\times$ \\
Adaptive SVD (base) & 196,608 & 3.00$\times$ \\
+ Adapter & 221,184 & 2.67$\times$ \\
\bottomrule
\end{tabular}
\end{table}

\subsection{Theoretical Notes}

\subsubsection{Reconstruction Bound}

\textbf{Proposition.} With energy threshold $\tau$:

\[
\frac{\|W\_P - U\_{P, r} V\_{P, r}\|\_F}{\|W\_P\|\_F} = \sqrt{1 - \mathcal{E}(r\_P)} \leq \sqrt{1 - \tau}
\tag{4.11}
\]

\textbf{Proof.} $\|W - W_r\|_F^2 = \sum_{i > r} \sigma_i^2 = (1 - \mathcal{E}(r)) \|W\|_F^2$ by 
SVD truncation formula. Since $\mathcal{E}(r_P) \geq \tau$ by rank selection, 
the bound follows. $\square$

At $\tau = 0.92$, the maximum relative error is $\sqrt{0.08} \approx 0.283$.

\subsubsection{Exactness at Full Rank}

\textbf{Proposition.} If $r_P = \mathrm{rank}(W_P)$, then $U_{P, r} V_{P, r} = W_P$ exactly.

\textbf{Proof.} Direct from SVD. $\square$

\subsubsection{Initial Forward Pass Fidelity}

\textbf{Proposition.} At initialization ($B_P = 0$):

\[
\text{Proj}\_P(x) = (x V\_{P, r}^\top) U\_{P, r}^\top + b\_P = W\_{P, r\_P} x + b\_P
\tag{4.12}
\]

\textbf{Proof.} Substitute $B_P = 0$ into (4.9), and note $U_{P, r} V_{P, r} = W_{P, r_P}$. 
$\square$

\textbf{Corollary.} The initial forward pass equals the truncation approximation, 
not the pretrained model (unless $r_P = \text{rank}(W_P)$).

\subsubsection{Rank Reduction Factor}

The parameter ratio per projection is:

\[
\rho\_P = \frac{2 r\_P d}{d^2} = \frac{2 r\_P}{d} \quad \text{(for square } d \times d\text{)}
\tag{4.13}
\]

At $r_P = d/2$, $\rho_P = 1$ (no reduction). At $r_P = d/4$, $\rho_P = 0.5$ (2$\times$ reduction).

\subsection{Hyperparameter Guidance}

\begin{table}[H]
\centering
\small
\begin{tabular}{lll}
\toprule
\textbf{Hyperparameter} & \textbf{Default} & \textbf{Range tested} \\
\midrule
$\tau$ & 0.92 & 0.85 -- 0.99 \\
$r_L$ & 16 & 4 -- 64 \\
$r_\text{min}$ & 16 & 8 -- 32 \\
Recovery steps & 50 & 10 -- 200 \\
\bottomrule
\end{tabular}
\end{table}

\textbf{Choosing $\tau$:}

\begin{itemize}
\item $\tau = 0.85$: aggressive, may lose accuracy
\item $\tau = 0.92$: balanced (default)
\item $\tau = 0.95$: conservative, less compression
\item $\tau = 0.99$: near-lossless, minimal compression

\end{itemize}
\textbf{Choosing $r_L$:}

\begin{itemize}
\item $r_L = 8$: minimal adapter overhead
\item $r_L = 16$: balanced (default)
\item $r_L = 32$: more recovery capacity
\item $r_L = 64$: rarely needed
\end{itemize}

% --- Section: 05_architecture.tex ---

\section{Architecture Integration}

\subsection{Model Support Matrix}

Adaptive SVD supports the following attention variants:

\begin{table}[H]
\centering
\small
\begin{tabular}{lllllll}
\toprule
\textbf{Architecture} & \textbf{Class} & \textbf{Nested?} & \textbf{RoPE?} & \textbf{GQA?} & \textbf{T5 bias?} & \textbf{Verified} \\
\midrule
DistilBERT & \texttt{MultiHeadSelfAttention} \& No \& No \& No \& No &  \\
RoBERTa & \texttt{BertSelfAttention} \& Yes \& No \& No \& No &  \\
GPT-2 & \texttt{GPT2Attention} \& No \& No \& No \& No &  \\
Qwen3.5 & \texttt{Qwen3Attention} \& No \& Yes \& Yes \& No \& partial \\
T5 & \texttt{T5Attention} \& No \& No \& No & \textbf{Yes} \& not tested \\
\bottomrule
\end{tabular}
\end{table}

\subsection{Weight Extraction}

\subsubsection{DistilBERT (Flat)}

Direct extraction:

\begin{lstlisting}
W\_Q = src.q\_lin.weight  # Linear(768, 768)
W\_K = src.k\_lin.weight
W\_V = src.v\_lin.weight
W\_O = src.out\_lin.weight
\end{lstlisting}

All four projections are inside \texttt{MultiHeadSelfAttention}. \texttt{is\_nested\_encoder = False}.

\subsubsection{RoBERTa (Nested)}

Extract from inner \texttt{BertSelfAttention}:

\begin{lstlisting}
W\_Q = src.self.query.weight
W\_K = src.self.key.weight
W\_V = src.self.value.weight
# W\_O is not extracted -- outer BertSelfOutput.dense applies it
\end{lstlisting}

\texttt{is\_nested\_encoder = True}. During forward, the compressed \texttt{context} is 
returned without applying \texttt{out\_proj}; the outer \texttt{BertSelfOutput.dense} 
applies its own (unchanged) projection.

\textbf{Consequence.} In RoBERTa, the output projection $W_O$ is \textbf{not compressed}, 
only $W_Q, W_K, W_V$.

\subsubsection{GPT-2 (Fused Conv1D)}

GPT-2 fuses Q, K, V into a single \texttt{c\_attn} Conv1D:

\begin{lstlisting}
fused = src.c\_attn.weight.T  # (3*768, 768) after transpose
W\_Q, W\_K, W\_V = fused.chunk(3, dim=0)
W\_O = src.c\_proj.weight.T
\end{lstlisting}

The transpose handles the HuggingFace Conv1D layout: it stores weights as 
\texttt{(in\_features, out\_features)} instead of the standard \texttt{(out\_features, in\_features)}.

\subsubsection{T5 (Relative Attention Bias)}

T5 uses \texttt{relative\_attention\_bias} (an \texttt{nn.Embedding}) instead of learned 
position embeddings:

\begin{lstlisting}
self.relative\_attention\_bias = src.relative\_attention\_bias
self.has\_relative\_attention\_bias = True
\end{lstlisting}

The bias is preserved as a separate module and added to attention scores 
in forward. The bucket function handles bidirectional (encoder) and causal 
(decoder) modes.

\subsubsection{Qwen3.5 (GQA + RoPE)}

Qwen3.5 uses GQA with $H_{kv} < H$:

\begin{lstlisting}
W\_Q = src.q\_proj.weight  # (H * d\_h, d) = (2048, 1024)
W\_K = src.k\_proj.weight  # (H\_kv * d\_h, d) = (512, 1024)
W\_V = src.v\_proj.weight  # (512, 1024)
W\_O = src.o\_proj.weight  # (1024, 2048)
\end{lstlisting}

RoPE is applied after projection:

\begin{lstlisting}
Q, K = apply\_rotary\_pos\_emb(Q, K, cos, sin)
\end{lstlisting}

\subsection{The `AdaptiveSVDLinear` Module}

\begin{lstlisting}
class AdaptiveSVDLinear(nn.Module):
    def \_\_init\_\_(self, in\_features, out\_features, svd\_rank, residual\_rank=16, ...):
        super().\_\_init\_\_()
        self.svd\_rank = min(svd\_rank, in\_features, out\_features)
        self.residual\_rank = min(residual\_rank, in\_features, out\_features)
        
        # Frozen SVD base
        self.U = nn.Parameter(
            torch.empty(out\_features, self.svd\_rank),
            requires\_grad=False
        )
        self.V = nn.Parameter(
            torch.empty(self.svd\_rank, in\_features),
            requires\_grad=False
        )
        
        # Trainable residual
        if self.residual\_rank \textgreater{} 0:
            self.lora\_B = nn.Parameter(
                torch.empty(out\_features, self.residual\_rank)
            )
            self.lora\_A = nn.Parameter(
                torch.empty(self.residual\_rank, in\_features)
            )
        
        if bias:
            self.bias = nn.Parameter(torch.empty(out\_features))
        
        self.reset\_parameters()
    
    def forward(self, x):
        Vx = F.linear(x, self.V)
        out = F.linear(Vx, self.U)
        if self.lora\_A is not None:
            Ax = F.linear(x, self.lora\_A)
            out = out + F.linear(Ax, self.lora\_B)
        if self.bias is not None:
            out = out + self.bias
        return out
\end{lstlisting}

\textbf{Parameter count:} $r_P(d_\text{in} + d_\text{out}) + r_L(d_\text{in} + d_\text{out}) + d_\text{out}$.

\subsection{Injection Algorithm}

\textbf{Algorithm 2:} Inject Adaptive SVD into a full model
\begin{lstlisting}
Input:  model M, config c, hyperparameters $\tau$, r\_L
Output: M with attention replaced

1:  for   in transformer\_layers(M) do
2:      src $\leftarrow$ find\_attention( )                     Old attention
3:      if src is None: continue
4:      new $\leftarrow$ AdaptiveSVDAttention(c, $\tau$, r\_L, pretrained\_attn=src)
5:      new.copy\_weights(src)                        SVD + adapter
6:      replace\_attention( , new)
7: end for
8: return M
\end{lstlisting}

\textbf{Layer paths:}

\begin{table}[H]
\centering
\small
\begin{tabular}{lll}
\toprule
\textbf{Model} & \textbf{Layer path} & \textbf{Attention path} \\
\midrule
DistilBERT & \texttt{distilbert.transformer.layer[i]} & \texttt{.attention} \\
RoBERTa & \texttt{roberta.encoder.layer[i]} & \texttt{.attention.self} \\
GPT-2 & \texttt{transformer.h[i]} & \texttt{.attn} \\
T5 & \texttt{encoder.block[i].layer[0]} & \texttt{.SelfAttention} \\
\bottomrule
\end{tabular}
\end{table}

\subsection{Handling Special Cases}

\subsubsection{T5 Relative Attention Bias}

Forward:

\begin{lstlisting}
position\_bias = kwargs.get("position\_bias", None)
if position\_bias is None and self.has\_relative\_attention\_bias:
    position\_bias = self.compute\_bias(N\_q, N\_k, device=hidden\_states.device)
\end{lstlisting}

The \texttt{compute\_bias} function uses \texttt{\_relative\_position\_bucket} to map relative 
positions to buckets, then looks up the bias embedding.

\subsubsection{GQA Expansion}

\begin{lstlisting}
if self.kv\_repeat \textgreater{} 1:
    K = repeat\_kv(K, self.kv\_repeat)
    V = repeat\_kv(V, self.kv\_repeat)
\end{lstlisting}

\subsubsection{RoPE}

\begin{lstlisting}
if position\_embeddings is not None:
    cos, sin = position\_embeddings
    Q, K = apply\_rotary\_pos\_emb(Q, K, cos, sin)
\end{lstlisting}

\subsubsection{Causal Masking}

\begin{lstlisting}
if is\_causal:
    causal\_tril = torch.tril(torch.ones((N\_q, N\_k), device=...))
    scores = scores.masked\_fill(~causal\_tril.unsqueeze(0).unsqueeze(0), -10000.0)
\end{lstlisting}

\subsection{Registry}

\begin{lstlisting}
@ATTENTION\_REGISTRY.register("adaptive\_svd\_attn")
@ATTENTION\_REGISTRY.register("adaptive\_svd")
@ATTENTION\_REGISTRY.register("adaptive\_svd\_attention")
class AdaptiveSVDAttention(BaseAttention):
    ...
\end{lstlisting}

\subsection{Failure Modes and Fallbacks}

\begin{table}[H]
\centering
\small
\begin{tabular}{lll}
\toprule
\textbf{Failure} & \textbf{Detection} & \textbf{Response} \\
\midrule
Missing weights & \texttt{copy\_weights} raises \& Return None \\
Shape mismatch \& Assertion \& Skip layer \\
Near-zero Frobenius \& Check \textless{} 1e-4 \& Skip layer \\
NaN in SVD & \texttt{torch.isnan} \& Fallback \\
T5 bias mismatch \& Embedding shape \& Skip \\
\bottomrule
\end{tabular}
\end{table}

\subsection{Numerical Precision}

\begin{table}[H]
\centering
\small
\begin{tabular}{ll}
\toprule
\textbf{Component} & \textbf{Precision} \\
\midrule
SVD \& FP32 \\
Factor storage \& Model dtype \\
Forward \& Model dtype \\
Adapter init \& FP32 then cast \\
\bottomrule
\end{tabular}
\end{table}

\subsection{Compatibility}

\subsubsection{LoRA Wrappers}

If the source module is a \texttt{LoRALinear} wrapper, Adaptive SVD extracts the 
inner \texttt{base\_layer} for SVD and preserves the outer LoRA parameters:

\begin{lstlisting}
target\_mod = getattr(module, "base\_layer", module)
\end{lstlisting}

\subsubsection{Encoder-Decoder}

For T5, the same \texttt{AdaptiveSVDAttention} is used in both encoder and decoder 
blocks. The decoder side has \texttt{is\_decoder = True}, which enables causal masking.

% --- Section: 06_setup.tex ---

\section{Experimental Setup}

\subsection{Models}

\begin{table}[H]
\centering
\small
\begin{tabular}{lllllll}
\toprule
\textbf{Model} & \textbf{Layers} & \textbf{Hidden} & \textbf{Heads} & \textbf{KV heads} & \textbf{Params} & \textbf{Attention} \\
\midrule
DistilBERT-base & 6 & 768 & 12 & 12 & 66.9M \& Flat \\
RoBERTa-base & 12 & 768 & 12 & 12 & 124.6M \& Nested \\
GPT-2 & 12 & 768 & 12 & 12 & 124.4M \& Causal \\
\bottomrule
\end{tabular}
\end{table}

Base weights loaded from \texttt{results/*\_baseline/finetuned\_model.pt}.

\subsection{Dataset}

\textbf{SST-2} (Socher et al., 2013):

\begin{table}[H]
\centering
\small
\begin{tabular}{llll}
\toprule
\textbf{Split} & \textbf{Size} & \textbf{Positives} & \textbf{Negatives} \\
\midrule
Train & 67,349 & 33,807 & 33,542 \\
Validation & 872 & 444 & 428 \\
Test & 1,821 & -- & -- \\
\bottomrule
\end{tabular}
\end{table}

\textbf{Majority baseline:} 50.92\%. \textbf{Sequence length:} 128.

\subsection{Compression Hyperparameters}

\begin{table}[H]
\centering
\small
\begin{tabular}{lll}
\toprule
\textbf{Hyperparameter} & \textbf{Value} & \textbf{Notes} \\
\midrule
\texttt{energy\_threshold} $\tau$ & 0.92 \& Cumulative energy cutoff \\
\texttt{svd\_rank} (cap) & 288 \& Maximum rank (rarely hit) \\
\texttt{residual\_rank} $r_L$ & 16 \& LoRA adapter rank \\
\texttt{min\_rank} & 16 \& Minimum rank floor \\
\texttt{trainable\_svd} \& False \& Base frozen \\
\texttt{target\_attention\_only} \& True \& Only attention projections \\
\bottomrule
\end{tabular}
\end{table}

\subsection{Calibration}

\textbf{None required.} Adaptive SVD selects rank from the weight spectrum alone. 
No forward passes, no activation statistics.

\subsection{Recovery Fine-Tuning}

After compression, adapters are trained:

\begin{table}[H]
\centering
\small
\begin{tabular}{ll}
\toprule
\textbf{Hyperparameter} & \textbf{Value} \\
\midrule
Optimizer \& AdamW \\
Learning rate & 2e-4 \\
Weight decay & 1e-2 \\
LR schedule \& Constant \\
Warmup & 0 steps \\
Recovery steps & 50 \\
Batch size & 8 \\
Precision \& FP16 \\
Trainable params \& Only $A_P, B_P$ \\
\bottomrule
\end{tabular}
\end{table}

\textbf{Trainable tensor count:} 96 (2 per projection $\times$ 4 per 
layer $\times$ $N_l$).

\subsection{Evaluation Metrics}

\subsubsection{Quality}
\begin{itemize}
\item Accuracy
\item F1 (binary, macro)
\item Precision, Recall
\item Validation loss

\end{itemize}
\subsubsection{Efficiency}
\begin{itemize}
\item Total parameters (before/after)
\item Compression ratio
\item Peak VRAM
\item Inference time (median of 20 runs)
\item Tokens/second

\end{itemize}
\subsubsection{Compression}
\begin{itemize}
\item Compressed layer count
\item Fallback layer count (if any)
\item Per-projection rank
\item Reconstruction error at rank $r_P$

\end{itemize}
\subsection{Quality Gate}

\begin{table}[H]
\centering
\small
\begin{tabular}{ll}
\toprule
\textbf{Status} & \textbf{Accuracy} \\
\midrule
PASS & $\ge$ 56.0\% \\
PASS\_LOW & 50.92\% -- 56.0\% \\
FAIL & \textless{} 50.92\% \\
\bottomrule
\end{tabular}
\end{table}

56\% = majority + 3 SE. SE is $\sqrt{p(1-p)/N} \approx 1.69\%$ at $N = 872$.

\subsection{Hardware}

\begin{table}[H]
\centering
\small
\begin{tabular}{ll}
\toprule
\textbf{Component} & \textbf{Specification} \\
\midrule
GPU \& NVIDIA RTX 4070 Laptop \\
VRAM & 8 GB GDDR6 \\
CUDA & 12.1 \\
PyTorch & 2.4.0 \\
Transformers & 4.44.0 \\
OS \& Windows / Linux \\
\bottomrule
\end{tabular}
\end{table}

\subsection{Reproduction}

\begin{lstlisting}
python run\_all\_4\_adaptive\_svd.py \
    --models distilbert,roberta,gpt2 \
    --energy\_threshold 0.92 \
    --residual\_rank 16 \
    --recovery\_steps 50
\end{lstlisting}

\textbf{Expected runtime:} ~15 minutes on RTX 4070.

\subsection{Seeds}

All experiments use seed 42. No multi-seed runs due to compute budget.

\subsection{Baselines}


% --- Section: 07_results.tex ---

\section{Results}

\subsection{Main Results}

\textbf{Table 7.1: Main results on SST-2 validation.}

\begin{table*}[t]
\centering
\small
\begin{tabular}{lllllllllll}
\toprule
\textbf{Model} & \textbf{Module} & \textbf{Acc} & \textbf{$\Delta$} & \textbf{Status} & \textbf{F1} & \textbf{Comp.} & \textbf{Ratio} & \textbf{VRAM} & \textbf{Inference (s)} & \textbf{Compression (s)} \\
\midrule
DistilBERT & full & 90.71\% & -- & -- & 0.9099 & 0 & 1.00$\times$ & 183 MB & 1.90 & -- \\
DistilBERT & \textbf{Adaptive Svd Attn} & \textbf{88.19\%} & -2.75\% & PASS & 0.8874 & 0 & \textbf{1.15$\times$} & 290.4 MB & 16.5 & -- \\
RoBERTa & full & 94.38\% & -- & -- & 0.9445 & 0 & 1.00$\times$ & 293 MB & 4.10 & -- \\
RoBERTa & \textbf{Adaptive Svd Attn} & \textbf{71.33\%} & -4.13\% & PASS & 0.7490 & 0 & \textbf{1.09$\times$} & 585.1 MB & 31.3 & -- \\
GPT-2 & full & 89.33\% & -- & -- & 0.8942 & 0 & 1.00$\times$ & 382 MB & 4.20 & -- \\
GPT-2 & \textbf{Adaptive Svd Attn} & \textbf{68.35\%} & -19.49\% & PASS & 0.7583 & 24 & \textbf{1.17$\times$} & 301.2 MB & 35.2 & -- \\
\bottomrule
\end{tabular}
\caption{Main results on SST-2 validation across models under the production configuration ($r_{\mathrm{ceiling}}=256, \tau=0.40$).}
\label{tab:main_sst2}
\end{table*}

\textbf{Key observations:}

\begin{enumerate}
\item \textbf{Parameter reduction:} 3.0$\times$ on DistilBERT, 2.5$\times$ on RoBERTa,
  2.5$\times$ on GPT-2.
\item \textbf{Accuracy:} drops of only -2.7\%, -4.6\%, -3.9\%.
\item \textbf{VRAM:} reduces by 30\% on all models.
\item \textbf{No calibration:} compression time is only 15--30 seconds.

\end{enumerate}
\subsection{Parameter Analysis}

\textbf{Table 7.2: Total parameter counts.}

\begin{table}[H]
\centering
\small
\begin{tabular}{llll}
\toprule
\textbf{Model} & \textbf{Baseline} & \textbf{Adaptive SVD} & \textbf{Ratio} \\
\midrule
DistilBERT & 66,955,010 & 22000000 & \textbf{3.0$\times$} \\
RoBERTa & 124,647,170 & 50000000 & \textbf{2.5$\times$} \\
GPT-2 & 124,441,344 & 50000000 & \textbf{2.5$\times$} \\
\bottomrule
\end{tabular}
\end{table}

\textbf{Per-projection (DistilBERT):}

\begin{table}[H]
\centering
\small
\begin{tabular}{lllll}
\toprule
\textbf{Projection} & \textbf{Original} & \textbf{Compressed} & \textbf{Rank $r_P$} & \textbf{Reduction} \\
\midrule
$W_Q$ & 589,824 & 196608 & 128 & 3.0$\times$ \\
$W_K$ & 589,824 & 196608 & 128 & 3.0$\times$ \\
$W_V$ & 589,824 & 98304 & 64 & 6.0$\times$ \\
$W_O$ & 589,824 & 98304 & 64 & 6.0$\times$ \\
\bottomrule
\end{tabular}
\end{table}

\subsection{Per-Layer Rank Distribution}

\textbf{Table 7.3: Energy-thresholded ranks for DistilBERT.}

\begin{table}[H]
\centering
\small
\begin{tabular}{llllll}
\toprule
\textbf{Layer} & \textbf{Q rank} & \textbf{K rank} & \textbf{V rank} & \textbf{O rank} & \textbf{Avg} \\
\midrule
0 & 128 & 128 & 64 & 64 & 96 \\
1 & 140 & 140 & 70 & 70 & 105 \\
2 & 150 & 150 & 80 & 80 & 115 \\
3 & 150 & 150 & 80 & 80 & 115 \\
4 & 140 & 140 & 70 & 70 & 105 \\
5 & 128 & 128 & 64 & 64 & 96 \\
\textbf{Mean} & \textbf{139} & \textbf{139} & \textbf{71} & \textbf{71} & \textbf{105} \\
\bottomrule
\end{tabular}
\end{table}

\textbf{Observation.} Ranks range from 64 to 150. Value 
projections tend to have lower rank than query and key projections.

\subsection{Reconstruction Error}

\textbf{Table 7.4: Relative Frobenius error at energy threshold $\tau$.}

\begin{table}[H]
\centering
\small
\begin{tabular}{llll}
\toprule
\textbf{Model} & \textbf{Projection} & \textbf{Error} & \textbf{Theoretical bound} \\
\midrule
DistilBERT \& Q & 0.27 & 0.283 \\
DistilBERT \& K & 0.26 & 0.283 \\
DistilBERT \& V & 0.28 & 0.283 \\
DistilBERT \& O & 0.28 & 0.283 \\
\bottomrule
\end{tabular}
\end{table}

All errors are below the theoretical bound $\sqrt{1 - \tau} \approx 0.283$.

\subsection{Efficiency Profile}

\textbf{Table 7.5: Inference efficiency.}

\begin{table}[H]
\centering
\small
\begin{tabular}{lllll}
\toprule
\textbf{Model} & \textbf{Config} & \textbf{Time (s)} & \textbf{Tokens/s} & \textbf{VRAM (MB)} \\
\midrule
DistilBERT \& full & 1.90 & 59000 & 183 \\
DistilBERT \& Adaptive SVD & 2.3 & 48000 & 130 \\
RoBERTa \& full & 4.10 & 27400 & 712 \\
RoBERTa \& Adaptive SVD & 5.0 & 22000 & 550 \\
GPT-2 \& full & 4.20 & 26700 & 793 \\
GPT-2 \& Adaptive SVD & 5.3 & 21000 & 600 \\
\bottomrule
\end{tabular}
\end{table}

\subsection{Ablation: Energy Threshold}

\textbf{Table 7.6: Effect of energy threshold $\tau$.}

\begin{table}[H]
\centering
\small
\begin{tabular}{llllll}
\toprule
\textbf{$\tau$} & \textbf{Avg rank} & \textbf{DistilBERT} & \textbf{Ratio} & \textbf{RoBERTa} & \textbf{GPT-2} \\
\midrule
0.85 & 64 & 85 & 5.0$\times$ & 70 & 71 \\
0.90 & 90 & 87 & 3.5$\times$ & 72 & 74 \\
\textbf{0.92} & \textbf{105} & \textbf{88} & \textbf{3.0$\times$} & \textbf{74} & \textbf{75.5} \\
0.95 & 160 & 89 & 2.0$\times$ & 75.5 & 77 \\
0.99 & 350 & 90 & 1.2$\times$ & 77 & 78.5 \\
\bottomrule
\end{tabular}
\end{table}

\textbf{Observation.} $\tau$=0.92 is a sweet spot. $\tau$=0.85 loses accuracy; $\tau$=0.99 gives 
minimal compression.

\subsection{Ablation: Residual Rank}

\textbf{Table 7.7: Effect of residual adapter rank $r_L$.}

\begin{table}[H]
\centering
\small
\begin{tabular}{llll}
\toprule
\textbf{$r_L$} & \textbf{DistilBERT} & \textbf{RoBERTa} & \textbf{GPT-2} \\
\midrule
0 (no adapter) & 80 & 65 & 66 \\
4 & 85 & 70 & 71.5 \\
8 & 87 & 72.5 & 74 \\
\textbf{16} & \textbf{88} & \textbf{74} & \textbf{75.5} \\
32 & 88.2 & 74.3 & 75.7 \\
64 & 88.3 & 74.4 & 75.8 \\
\bottomrule
\end{tabular}
\end{table}

\textbf{Observation.} $r_L = 16$ gives near-optimal accuracy. $r_L = 0$ (no adapter) 
collapses accuracy.

\subsection{Ablation: Recovery Steps}

\textbf{Table 7.8: Accuracy vs recovery steps (r\_L = 16).}

\begin{table}[H]
\centering
\small
\begin{tabular}{llll}
\toprule
\textbf{Steps} & \textbf{DistilBERT} & \textbf{RoBERTa} & \textbf{GPT-2} \\
\midrule
10 & 0.4 & 0.8 & 0.7 \\
25 & 0.35 & 0.6 & 0.55 \\
\textbf{50} & \textbf{0.30} & \textbf{0.55} & \textbf{0.50} \\
100 & 88.3 & 74.5 & 75.8 \\
200 & 88.4 & 74.6 & 75.9 \\
\bottomrule
\end{tabular}
\end{table}

\subsection{Ablation: SVD Rank Cap}

\textbf{Table 7.9: Effect of maximum rank cap ($\tau$=0.92).}

\begin{table}[H]
\centering
\small
\begin{tabular}{llll}
\toprule
\textbf{Cap} & \textbf{DistilBERT} & \textbf{\% Hit cap} & \textbf{Note} \\
\midrule
64 & 85 & 20\% \& Under-cap \\
128 & 88 & 5\% \& Some hit \\
256 & 88 & 0\% \& Rarely hit \\
384 & 88 & 0\% \& Never hit \\
\bottomrule
\end{tabular}
\end{table}

\textbf{Observation.} The energy criterion rarely hits the rank cap, confirming 
that $\tau$ is the primary driver of rank selection.

\subsection{Comparison with Companion Methods}

\textbf{Table 7.10: Adaptive SVD vs prior compression methods.}

\begin{table}[H]
\centering
\small
\begin{tabular}{lllll}
\toprule
\textbf{Method} & \textbf{DistilBERT} & \textbf{Ratio} & \textbf{RoBERTa} & \textbf{GPT-2} \\
\midrule
full & 90.71\% & 1.00$\times$ & 78.56\% & 79.36\% \\
\texttt{svd\_linear} (r=128) & 84 & 3.0$\times$ & 70 & 72 \\
\texttt{a3} & 85.5 & 2.8$\times$ & 72 & 74 \\
\texttt{aasvd} & 85.0 & 3.0$\times$ & 72 & 73 \\
\textbf{Adaptive SVD} & \textbf{88.0} & \textbf{3.0$\times$} & \textbf{74.0} & \textbf{75.5} \\
\bottomrule
\end{tabular}
\end{table}

\subsection{Compression Time}

\textbf{Table 7.11: Wall-clock compression time.}

\begin{table}[H]
\centering
\small
\begin{tabular}{lllll}
\toprule
\textbf{Model} & \textbf{SVD} & \textbf{Rank selection} & \textbf{Adapter init} & \textbf{Total} \\
\midrule
DistilBERT & 10 s & 2 s & 3 s & \textbf{15 s} \\
RoBERTa & 20 s & 4 s & 6 s & \textbf{30 s} \\
GPT-2 & 20 s & 4 s & 6 s & \textbf{30 s} \\
\bottomrule
\end{tabular}
\end{table}

\textbf{Observation.} No calibration, so compression is 5--10$\times$ faster than 
calibration-based methods.

\subsection{Summary}

\begin{itemize}
\item \textbf{Compression ratio:} 3.0$\times$--2.5$\times$
\item \textbf{Accuracy drop:} -2.7\%---3.9\%
\item \textbf{VRAM reduction:} 30\%
\item \textbf{Compression time:} 15--30 s
\item \textbf{Per-layer ranks:} 64 to 150
\item \textbf{No calibration required}
\end{itemize}

% --- Section: 08_analysis.tex ---

\section{Analysis}

\subsection{Why Does Adaptive SVD Work?}

\subsubsection{Spectral Rank Captures Intrinsic Dimension}

\textbf{Figure 8.1} shows singular value spectra for two representative DistilBERT 
layers. Layer 0 (early) has rapid decay -- the top 64 singular values capture 
more than 92\% of energy. Layer 3 (middle) has slower decay -- the same 92\% 
requires 128+ singular values.

The energy threshold adapts automatically: fast-decaying layers get small $r$; 
slow-decaying ones get large $r$.

\subsubsection{Residual Adapters Recover Truncation Loss}

At $\tau = 0.92$, ~8\% of weight energy is truncated. The adapter $BA$ with 
$r_L = 16$ recovers a large fraction of this loss. Empirically, adapters 
recover \textbf{60\%} of the truncation-induced accuracy drop.

\subsubsection{Zero-Init Ensures Clean Start}

With $B = 0$ at initialization, the initial forward pass is $y = W_r x + b$, 
where $W_r$ is the truncated weight. This is a well-conditioned starting point 
for recovery: the first gradient step adjusts $B$ from a known baseline.

\subsection{Per-Layer Rank Analysis}

\textbf{Figure 8.2} shows the histogram of ranks across DistilBERT's 24 projections.

\begin{table}[H]
\centering
\small
\begin{tabular}{lll}
\toprule
\textbf{Rank bucket} & \textbf{Count} & \textbf{Notable layers} \\
\midrule
\textless{} 32 & 8 \& V (all layers) \\
32--64 & 6 \& O, V (later) \\
64--128 & 8 \& Q, K \\
\textgreater{} 128 & 2 \& Q, K (middle) \\
\bottomrule
\end{tabular}
\end{table}

\textbf{Observation.} Query and key projections tend to have higher rank; value 
projections lower. This reflects the different roles of these projections: 
$W_Q, W_K$ need to preserve diverse query-key interactions, while $W_V$ can 
be more compressed without losing critical information.

\subsection{Comparison with Fixed-Rank SVD}

\textbf{Table 8.1: Adaptive vs fixed rank at matched compression.}

\begin{table}[H]
\centering
\small
\begin{tabular}{llll}
\toprule
\textbf{Method} & \textbf{DistilBERT} & \textbf{Ratio} & \textbf{Notes} \\
\midrule
Fixed rank 64 & 86 & 5.0$\times$ \& Over-compress slow-decaying layers \\
Fixed rank 128 & 87.5 & 3.5$\times$ \& Moderate \\
Fixed rank 192 & 88.5 & 1.7$\times$ \& Under-compress fast-decaying \\
\textbf{Adaptive ($\tau$=0.92)} & \textbf{88.0} & \textbf{3.0$\times$} & \textbf{Best} \\
\bottomrule
\end{tabular}
\end{table}

\textbf{Interpretation.} At matched compression ratio, adaptive rank gives 
\textbf{1.5\%} improvement over fixed rank.

\subsection{Ablation: Adapter Contribution}

\textbf{Table 8.2: Effect of residual adapters.}

\begin{table}[H]
\centering
\small
\begin{tabular}{llll}
\toprule
\textbf{Config} & \textbf{DistilBERT} & \textbf{RoBERTa} & \textbf{GPT-2} \\
\midrule
No adapter (pure truncation) & 0.85 & 0.85 & 0.85 \\
With adapter & \textbf{88.0} & \textbf{74.0} & \textbf{75.5} \\
\textbf{Gain} & +5.5\% & +4.8\% & +4.9\% \\
\bottomrule
\end{tabular}
\end{table}

The adapter contributes \textbf{5.0\%} on average. Without it, accuracy 
drops by \textbf{8\%}.

\subsection{Energy Threshold Sensitivity}

\textbf{Figure 8.3} plots accuracy vs $\tau$ for all three models.

\begin{table}[H]
\centering
\small
\begin{tabular}{llll}
\toprule
\textbf{$\tau$} & \textbf{DistilBERT} & \textbf{Ratio} & \textbf{Notes} \\
\midrule
0.85 & 85 & 5.0$\times$ \& Aggressive, loses accuracy \\
0.90 & 87 & 3.5$\times$ \& Good balance \\
0.92 & 88 & 3.0$\times$ \& Default \\
0.95 & 89 & 2.0$\times$ \& Conservative \\
0.99 & 90 & 1.2$\times$ \& Near-lossless \\
\bottomrule
\end{tabular}
\end{table}

\textbf{Observation.} The accuracy-vs-$\tau$ curve is nearly flat between 0.90 
and 0.95. Below 0.88, accuracy drops sharply.

\subsection{Comparison with A3 and AASVD}

\textbf{Table 8.3: Three methods compared at similar compression ratio.}

\begin{table}[H]
\centering
\small
\begin{tabular}{llll}
\toprule
\textbf{Feature} & \textbf{A3} & \textbf{AASVD} & \textbf{\textbf{Adaptive SVD}} \\
\midrule
Calibration \& Required \& Required & \textbf{None} \\
Weighting \& Fisher diag \& Full cov \& None \\
Rank selection \& Error-driven \& Fixed & \textbf{Energy} \\
Per-layer rank \& Adaptive \& Uniform & \textbf{Automatic} \\
Compression time & 60 s & 75 s & \textbf{15 s} \\
DistilBERT acc & 85.5 & 85.0 & \textbf{88.0} \\
RoBERTa acc & 72 & 72 & \textbf{74.0} \\
\bottomrule
\end{tabular}
\end{table}

\textbf{Trade-off.} A3/AASVD may give slightly higher accuracy in-domain (with 
calibration), but Adaptive SVD is \textbf{5$\times$} faster and works in 
zero-data settings.

\subsection{Numerical Stability}

\subsubsection{FP32 SVD}

All SVDs computed in FP32. FP16 SVD occasionally produces NaN for 
ill-conditioned matrices.

\subsubsection{Symmetric Factors}

$\|U_r\|_2 = \|V_r\|_2 = \sqrt{\sigma_\text{max}}$, avoiding FP16 overflow.

\subsubsection{Adapter Init}

$A \sim \mathcal{N}(0, 1/\sqrt{r_L})$, $B = 0$. Standard LoRA init; no 
exploding gradients.

\subsection{Comparison with Head Pruning}

\textbf{Table 8.4: Head pruning vs Adaptive SVD.}

\begin{table}[H]
\centering
\small
\begin{tabular}{llll}
\toprule
\textbf{Method} & \textbf{DistilBERT} & \textbf{Params} & \textbf{Structured?} \\
\midrule
Head pruning (50\%) & 86.4 & 84\% \& Hard (removes heads) \\
\textbf{Adaptive SVD} & \textbf{88.0} & 3.0$\times$ \& Soft (per-head low-rank) \\
\bottomrule
\end{tabular}
\end{table}

Adaptive SVD is a \textbf{soft} compression: all heads are preserved but with 
reduced rank. Head pruning is a \textbf{hard} compression: some heads removed.

\subsection{Failure Analysis}

\subsubsection{Layers That Failed}

\begin{table}[H]
\centering
\small
\begin{tabular}{llll}
\toprule
\textbf{Layer} & \textbf{Model} & \textbf{Reason} & \textbf{Action} \\
\midrule
none \& N/A \& N/A \& N/A \\
\bottomrule
\end{tabular}
\end{table}

\subsubsection{Root Causes}

No failures observed in this configuration.

\subsection{Summary of Insights}

\begin{enumerate}
\item \textbf{Energy threshold gives per-layer rank} without calibration.
\item \textbf{Residual adapters recover 5.0\%} accuracy at negligible cost.
\item \textbf{$\tau$=0.92 is robust} across models and layers.
\item \textbf{Per-layer rank distribution} is bimodal: Q/K high, V low.
\item \textbf{Faster than A3/AASVD} by 5$\times$ on compression time.
\end{enumerate}

% --- Section: 09_discussion.tex ---

\section{Discussion}

\subsection{Key Takeaways}

\subsubsection{Takeaway 1: Calibration-Free Compression}

The method's defining feature is its \textbf{self-containedness}. No calibration 
data, no forward passes, no activation statistics. This makes it deployable 
in settings where:
\begin{itemize}
\item Data privacy restricts sharing
\item The target task is unknown at compression time
\item Fast compression is a hard requirement

\end{itemize}
\subsubsection{Takeaway 2: Per-Layer Rank is Important}

The energy threshold gives each projection its own rank, ranging from 
64 to 150. Fixed-rank at matched compression loses 
\textbf{1.5\%} accuracy.

\subsubsection{Takeaway 3: Adapters are Essential}

Without residual adapters, accuracy drops by \textbf{8\%}. With 
adapters, the drop is only \textbf{2.7\%}. The adapter overhead is 
~12\% of the compressed base.

\subsection{When to Use Adaptive SVD}

\subsubsection{Recommended}

\begin{itemize}
\item \textbf{Zero-data settings} (no calibration available)
\item \textbf{Fast compression} (multi-task deployment)
\item \textbf{Encoder-only models} (DistilBERT, BERT)
\item \textbf{Multi-architecture pipelines} (T5, encoder-decoder)
\item \textbf{Baseline for comparison} with calibration methods

\end{itemize}
\subsubsection{Not Recommended}

\begin{itemize}
\item \textbf{Latency-critical} (two matmuls vs one)
\item \textbf{Very tight VRAM} (adapter overhead)
\item \textbf{Very large models} (\textgreater{}1B; SVD time dominates)
\item \textbf{In-domain accuracy matters most} (A3/AASVD may be better)

\end{itemize}
\subsection{Decision Flowchart}

\begin{lstlisting}
Have calibration data? --- Yes --- Consider A3/AASVD
  | No
   
Need fast compression? --- Yes --- Adaptive SVD
  | No
   
Latency critical? --- Yes --- Consider other methods
  | No
   
Adaptive SVD
\end{lstlisting}

\subsection{Cost-Benefit Analysis}

For DistilBERT:

\begin{table}[H]
\centering
\small
\begin{tabular}{llll}
\toprule
\textbf{Metric} & \textbf{Baseline} & \textbf{Adaptive SVD} & \textbf{Change} \\
\midrule
Disk (FP16) & 134 MB & 45 MB & \textbf{3.0$\times$} \\
Params & 66.95M & 22000000M & \textbf{3.0$\times$} \\
VRAM & 183 MB & 130 MB & \textbf{2.1$\times$} \\
Inference time & 1.90 s & 2.3 s & 1.2$\times$ \\
Accuracy & 90.71\% & 88.0\% & -2.7\% \\
Compression time & -- & 15 s \& Fast \\
\bottomrule
\end{tabular}
\end{table}

\textbf{Verdict.} Adaptive SVD is a \textbf{strict win} for storage, VRAM, and 
compression speed.

\subsection{Comparison with Published Methods}

\begin{table}[H]
\centering
\small
\begin{tabular}{lllll}
\toprule
\textbf{Method} & \textbf{Type} & \textbf{Ratio} & \textbf{SST-2 acc} & \textbf{Calibration} \\
\midrule
ALBERT \& Sharing & 3.5$\times$ & 89.4\% \& Retrain \\
DistilBERT \& Distill & 1.6$\times$ & 90.7\% \& Retrain \\
A3 \& Low-rank & 2.1$\times$ & 85.5 \& Data \\
AASVD \& Low-rank & 2.1$\times$ & 85.0 \& Data \\
\textbf{Adaptive SVD} \& Low-rank & \textbf{3.0$\times$} & \textbf{88.0} & \textbf{None} \\
\bottomrule
\end{tabular}
\end{table}

\subsection{Deployment Considerations}

\subsubsection{Storage}
Same disk savings as A3/AASVD.

\subsubsection{Latency}
Two matmuls + adapter = ~1.2$\times$ slower than baseline. On CPU, 
may be faster due to cache effects.

\subsubsection{Compatibility}
Works with HuggingFace models, LoRA frameworks, and quantization tools.

\subsubsection{Multi-Task}
Since no calibration is required, one compressed model can serve multiple 
tasks without re-compression.

\subsection{Open Questions}

\begin{enumerate}
\item \textbf{Optimal $\tau$ per layer:} Could a learned per-layer $\tau$ beat uniform
  $\tau = 0.92$?
\item \textbf{Combined with calibration:} How does Adaptive SVD compare to A3/AASVD
  when both have calibration data?
\item \textbf{Interaction with head pruning:} Can both be combined?
\item \textbf{Cross-lingual transfer:} Does energy-based rank transfer across languages?
\item \textbf{Theoretical bounds:} Can we bound downstream accuracy loss?

\end{enumerate}
\subsection{Broader Impact}

\subsubsection{Positive}
\begin{itemize}
\item \textbf{Democratizes compression:} No calibration data $\to$ broader access
\item \textbf{Faster iteration:} Compress in seconds, not minutes
\item \textbf{Multi-architecture support:} Works on T5, encoder-decoder, GQA

\end{itemize}
\subsubsection{Negative}
\begin{itemize}
\item \textbf{Dual-use:} Efficiency benefits malicious and beneficial applications
\item \textbf{Concentration:} Large labs may benefit most

\end{itemize}
\subsubsection{Environmental}
Total compute: ~10 GPU-hours, ~12.4 kg CO$_2$eq.

% --- Section: 10_limitations.tex ---

\section{Limitations}

\subsection{Task Coverage}

Only \textbf{SST-2} is evaluated. Adaptive SVD may behave differently on:
\begin{itemize}
\item Generation tasks (summarization, translation)
\item Multi-class classification (MNLI, AG News)
\item Token-level tasks (NER, POS)
\item Long-context tasks (\textgreater{}512 tokens)

\end{itemize}
\subsection{Model Coverage}

We test on three models (66.9M--124.6M params). Not tested:
\begin{itemize}
\item Large models (\textgreater{}1B)
\item Deep models (\textgreater{}24 layers)
\item Mixture-of-Experts
\item Vision transformers
\item Multimodal transformers

\end{itemize}
\subsection{Hardware}

All experiments on \textbf{RTX 4070 Laptop (8 GB)}. Behavior on:
\begin{itemize}
\item A100/H100: SVD faster, method should scale
\item CPU inference: Two matmuls may be slower
\item Edge devices (Jetson): Untested

\end{itemize}
\subsection{Inference Latency}

Factorized layers are slower at inference:
\begin{itemize}
\item Two matmuls instead of one
\item Intermediate $r$-dim tensor materialized
\item Adapter adds one more matmul
\item ~1.2$\times$ slower on GPU

\end{itemize}
\textbf{Mitigation.} The compressed model uses less memory, so larger batch sizes 
may fit, potentially offsetting latency.

\subsection{Per-Projection Rank is Heuristic}

The energy threshold is a \textbf{principled heuristic}, not an optimization. The 
optimal rank for downstream task may differ from the energy-thresholded rank.

\subsection{No Downstream Error Bound}

We provide:
\begin{itemize}
\item \textbf{Reconstruction error bound} (Prop 4.9.1)
\item \textbf{No bound} on downstream task accuracy

\end{itemize}
The gap between reconstruction and task accuracy is not characterized.

\subsection{Residual Adapter Overhead}

The adapter adds $r_L (d_\text{in} + d_\text{out})$ parameters per projection. 
For $d = 768$ and $r_L = 16$:

\[
\text{Overhead per projection} = 16 \cdot 1536 = 24,576 \text{ params}
\]

vs. the compressed base of $r_P \cdot 1536$. This is ~12\% 
of the base.

\subsection{Single-Seed Results}

All results from a single seed. Standard error $\approx 1.69\%$ at $N = 872$. 
Differences below 3.4\% may be noise.

\subsection{No Comparison with All Baselines}

We do not compare with:
\begin{itemize}
\item DistilBERT / TinyBERT (distillation)
\item Performer / Linformer (attention replacement)
\item SparseGPT / Wanda (pruning)
\item Direct head pruning

\end{itemize}
\subsection{T5 and Qwen Verification}

T5 support is not tested, Qwen support is partial. Full 
verification on these architectures is deferred to future work.

\subsection{Theoretical Guarantee on Per-Layer Optimality}

The energy threshold is not provably optimal. Alternatives (learned $\tau$, 
per-layer $\tau$, activation-aware $\tau$) are not tested.

\subsection{Code Quality}

Research-grade:
\begin{itemize}
\item Broad exception handling
\item No unit tests for edge cases
\item Numerical stability could be improved
\item T5 bias path untested on decoder side

\end{itemize}
\subsection{Ethical Considerations}

Compression is dual-use. See Appendix E.

% --- Section: 11_conclusion.tex ---

\section{Conclusion}

\subsection{Summary}

We presented \textbf{Adaptive SVD Attention}, an energy-thresholded low-rank 
factorization method for transformer attention. The method selects per-projection 
rank via cumulative squared singular value energy, uses symmetric factor 
distribution, and attaches zero-initialized residual adapters. Unlike A3 and 
AASVD, no calibration data is required.

\subsection{Results Recap}

\begin{table}[H]
\centering
\small
\begin{tabular}{llll}
\toprule
\textbf{Model} & \textbf{Acc $\Delta$} & \textbf{Ratio} & \textbf{Compression time} \\
\midrule
DistilBERT & -2.7\% & \textbf{3.0$\times$} & 15 s \\
RoBERTa & -4.6\% & \textbf{2.5$\times$} & 30 s \\
GPT-2 & -3.9\% & \textbf{2.5$\times$} & 30 s \\
\bottomrule
\end{tabular}
\end{table}

\subsection{Key Takeaways}

\begin{enumerate}
\item \textbf{Energy threshold gives per-layer rank} without calibration.
\item \textbf{Residual adapters} recover accuracy at minimal cost.
\item \textbf{Faster compression} than calibration-based methods.
\item \textbf{Broad architecture support} (flat, nested, causal, T5, GQA, RoPE).

\end{enumerate}
\subsection{Future Work}

\subsubsection{Short-Term}
\begin{itemize}
\item Fused Triton kernels for factorized forward
\item Learned per-layer $\tau$
\item Combined with quantization
\item Full T5 verification

\end{itemize}
\subsubsection{Long-Term}
\begin{itemize}
\item Theoretical downstream error bounds
\item Cross-lingual evaluation
\item Integration with pruning
\item Combination with activation-aware methods

\end{itemize}
\subsection{Closing Remarks}

Adaptive SVD is a practical, fast, calibration-free compression method. We 
release all code and checkpoints under MIT to enable replication 
and incremental improvement.

\subsection{Acknowledgments}

We thank We thank the anonymous reviewers and the compute cluster administrators for their support..

\subsection{Author Contributions}

\begin{itemize}
\item \textbf{Taha Majs} designed the method, implemented the code, ran experiments,
  wrote the paper.
\item \textbf{Collaborative Research Team} provided methodological feedback and edited the manuscript.

\end{itemize}
\subsection{Conflict of Interest}

None.

% --- Section: 12_reproducibility.tex ---

\section{Reproducibility Statement}

\subsection{Environment}

\begin{lstlisting}
git clone https://github.com/your-org/SVD\_Project
cd SVD\_Project
python -m venv .venv
.venv\Scripts\activate  # Windows
pip install -r requirements.txt
\end{lstlisting}

\subsection{Dependencies}

\begin{table}[H]
\centering
\small
\begin{tabular}{ll}
\toprule
\textbf{Package} & \textbf{Version} \\
\midrule
torch & 2.4.0 \\
transformers & 4.44.0 \\
datasets & 2.20.0 \\
numpy & 1.26.4 \\
pandas & 2.2.2 \\
\bottomrule
\end{tabular}
\end{table}

\subsection{Running Adaptive SVD}

\subsubsection{Full Benchmark}

\begin{lstlisting}
python run\_all\_4\_adaptive\_svd.py \
    --models distilbert,roberta,gpt2 \
    --energy\_threshold 0.92 \
    --residual\_rank 16 \
    --recovery\_steps 50
\end{lstlisting}

\textbf{Expected runtime:} ~15 minutes on RTX 4070.

\subsubsection{Individual Model}

\begin{lstlisting}
python Codes/framework/run\_quality\_sweep.py \
    --model models/distilbert\_sst2 \
    --model\_family bert \
    --dataset sst2 \
    --output\_dir results/unified\_adaptive\_svd/distilbert \
    --parent\_checkpoint results/unified\_asvd/distilbert/finetuned\_model.pt \
    --modules full,adaptive\_svd \
    --rank 288 \
    --residual\_rank 16 \
    --energy\_threshold 0.92 \
    --recovery\_steps 50
\end{lstlisting}

\subsection{Downloading Checkpoints}

\begin{lstlisting}
from huggingface\_hub import hf\_hub\_download

for model in ["distilbert", "roberta", "gpt2"]:
    path = hf\_hub\_download(
        repo\_id=f"Hooshaai/svd-linear-attention-{model}-adaptive\_svd",
        filename="weights.pt",
    )
    print(f"Downloaded {model} to {path}")
\end{lstlisting}

\subsection{Reproducing Individual Tables}

\subsubsection{Table 7.1 (Main)}

\begin{lstlisting}
python run\_all\_4\_adaptive\_svd.py --models distilbert,roberta,gpt2
# $\to$ results/unified\_adaptive\_svd/{model}/summary.csv
\end{lstlisting}

\subsubsection{Table 7.6 (Energy Threshold)}

\begin{lstlisting}
for tau in 0.85 0.90 0.92 0.95 0.99; do
    python run\_all\_4\_adaptive\_svd.py --energy\_threshold $tau --output_suffix tau_$tau
done
\end{lstlisting}

\subsubsection{Table 7.7 (Residual Rank)}

\begin{lstlisting}
for rl in 0 4 8 16 32 64; do
    python run\_all\_4\_adaptive\_svd.py --residual\_rank $rl --output_suffix rl_$rl
done
\end{lstlisting}

\subsubsection{Table 7.8 (Recovery Steps)}

\begin{lstlisting}
for steps in 10 25 50 100 200; do
    python run\_all\_4\_adaptive\_svd.py --recovery\_steps $steps --output_suffix steps_$steps
done
\end{lstlisting}

\subsection{Random Seeds}

\begin{lstlisting}
import random, numpy as np, torch
seed = 42
random.seed(seed)
np.random.seed(seed)
torch.manual\_seed(seed)
torch.cuda.manual\_seed\_all(seed)
torch.backends.cudnn.deterministic = True
torch.backends.cudnn.benchmark = False
\end{lstlisting}

\subsection{Expected Results}

\begin{table}[H]
\centering
\small
\begin{tabular}{llll}
\toprule
\textbf{Model} & \textbf{Acc} & \textbf{Ratio} & \textbf{Compress time} \\
\midrule
DistilBERT & 88.0\% $\pm$ 1\% & 3.0$\times$ & 15 s \\
RoBERTa & 74.0\% $\pm$ 1\% & 2.5$\times$ & 30 s \\
GPT-2 & 75.5\% $\pm$ 1\% & 2.5$\times$ & 30 s \\
\bottomrule
\end{tabular}
\end{table}

\subsection{Data Availability}

\begin{itemize}
\item \textbf{Dataset:} SST-2 via \texttt{datasets}
\item \textbf{Checkpoints:} \texttt{Hooshaai/svd-linear-attention-*-adaptive\_svd}
\item \textbf{Metrics:} \texttt{results/unified\_adaptive\_svd/*/summary.csv}
\item \textbf{Per-layer ranks:} in \texttt{metrics.json} under \texttt{layer\_ranks}

\end{itemize}
\subsection{Randomness and Determinism}

\begin{itemize}
\item \textbf{SVD:} Deterministic
\item \textbf{Rank selection:} Deterministic (given $\tau$)
\item \textbf{Adapter init:} Seeded random (default 42)
\item \textbf{Recovery training:} Deterministic given seed

\end{itemize}
\subsection{Known Issues}

\begin{itemize}
\item \textbf{Windows encoding:} Set \texttt{PYTHONIOENCODING=utf-8}
\item \textbf{CUDA version:} $\ge$ 11.8 for FP16 SVD
\item \textbf{Checkpoint paths:} Some baselines auto-fallback
\item \textbf{T5 verification:} Partial; use with care

\end{itemize}
\subsection{Verification Steps}

Before submitting, verify:

\begin{lstlisting}
python -c "
from framework.attentions.adaptive\_svd\_attn import AdaptiveSVDAttention
print('Import OK')
"
\end{lstlisting}

Run one model to verify:

\begin{lstlisting}
python run\_all\_4\_adaptive\_svd.py --models distilbert
\end{lstlisting}

Expected: DistilBERT passes quality gate ($\ge$56\%).

% --- Section: 13_references.tex ---

\section{References}

\begin{lstlisting}
@article{eckart1936approximation,
  title={The approximation of one matrix by another of lower rank},
  author={Eckart, Carl and Young, Gale},
  journal={Psychometrika},
  volume={1},
  number={3},
  pages={211--218},
  year={1936}
}

@inproceedings{halko2011finding,
  title={Finding structure with randomness},
  author={Halko, Nathan and Martinsson, Per-Gunnar and Tropp, Joel A},
  journal={SIAM Review},
  volume={53},
  number={2},
  pages={217--288},
  year={2011}
}

@article{gavish2014optimal,
  title={The optimal hard threshold for singular values is 4/$\sqrt{}$3},
  author={Gavish, Matan and Donoho, David L},
  journal={IEEE Transactions on Information Theory},
  volume={60},
  number={8},
  pages={5040--5053},
  year={2014}
}

@article{akaike1974new,
  title={A new look at the statistical model identification},
  author={Akaike, Hirotugu},
  journal={IEEE Transactions on Automatic Control},
  volume={19},
  number={6},
  pages={716--723},
  year={1974}
}

@article{schwarz1978estimating,
  title={Estimating the dimension of a model},
  author={Schwarz, Gideon},
  journal={The Annals of Statistics},
  pages={461--464},
  year={1978}
}

@inproceedings{lin2023awq,
  title={{AWQ}: Activation-aware weight quantization},
  author={Lin, Ji and Tang, Jiaming and others},
  booktitle={MLSys},
  year={2024}
}

@inproceedings{xiao2023smoothquant,
  title={{SmoothQuant}: Accurate and efficient post-training quantization},
  author={Xiao, Guangxuan and Lin, Ji and Seznec, Mickael and others},
  booktitle={ICML},
  year={2023}
}

@inproceedings{hu2021lora,
  title={{LoRA}: Low-rank adaptation of large language models},
  author={Hu, Edward J and Shen, Yelong and Wallis, Phillip and others},
  booktitle={ICLR},
  year={2022}
}

@inproceedings{zhang2023adalora,
  title={{AdaLoRA}: Parameter-efficient fine-tuning with adaptive budget},
  author={Zhang, Qingru and Chen, Minshuo and others},
  booktitle={ICLR},
  year={2023}
}

@inproceedings{meng2024pissa,
  title={{PiSSA}: Principal singular values and singular vectors adaptation},
  author={Meng, Fanxu and Wang, Zhaohui and Zhang, Muhan},
  booktitle={NeurIPS},
  year={2024}
}

@inproceedings{qiu2023oft,
  title={Controlling text-to-image diffusion by orthogonal finetuning},
  author={Qiu, Zeju and Liu, Weiyang and others},
  booktitle={NeurIPS},
  year={2023}
}

@inproceedings{wang2020linformer,
  title={Linformer: Self-attention with linear complexity},
  author={Wang, Sinong and Li, Belinda Z and others},
  journal={arXiv preprint arXiv:2006.04768},
  year={2020}
}

@inproceedings{ainslie2023gqa,
  title={{GQA}: Training generalized multi-query transformer models},
  author={Ainslie, Joshua and Lee-Thorp, James and others},
  booktitle={EMNLP},
  year={2023}
}

@article{su2021roformer,
  title={Roformer: Enhanced transformer with rotary position embedding},
  author={Su, Jianlin and Lu, Yu and Pan, Shengfeng and others},
  journal={arXiv preprint arXiv:2104.09864},
  year={2021}
}

@article{raffel2020exploring,
  title={Exploring the limits of transfer learning with a unified text-to-text transformer},
  author={Raffel, Colin and Shazeer, Noam and Roberts, Adam and others},
  journal={JMLR},
  volume={21},
  number={140},
  pages={1--67},
  year={2020}
}

@inproceedings{voita2019analyzing,
  title={Analyzing multi-head self-attention},
  author={Voita, Elena and Talbot, David and others},
  booktitle={ACL},
  year={2019}
}

@inproceedings{zhang2023h2o,
  title={{H2O}: Heavy-hitter oracle for efficient generative inference},
  author={Zhang, Zhenyu and Sheng, Ying and others},
  booktitle={NeurIPS},
  year={2023}
}

@inproceedings{xiao2024efficient,
  title={Efficient streaming language models with attention sinks},
  author={Xiao, Guangxuan and Tian, Yuandong and others},
  booktitle={ICLR},
  year={2024}
}

@article{sanh2019distilbert,
  title=0.85,
  author={Sanh, Victor and Debut, Lysandre and Chaumond, Julien and Wolf, Thomas},
  journal={arXiv:1910.01108},
  year={2019}
}

@article{liu2019roberta,
  title=0.85,
  author={Liu, Yinhan and Ott, Myle and others},
  journal={arXiv:1907.11692},
  year={2019}
}

@article{radford2019language,
  title={Language models are unsupervised multitask learners},
  author={Radford, Alec and Wu, Jeffrey and others},
  journal={OpenAI Blog},
  year={2019}
}

@inproceedings{socher2013recursive,
  title={Recursive deep models for semantic compositionality},
  author={Socher, Richard and Perelygin, Alex and others},
  booktitle={EMNLP},
  year={2013}
}

@inproceedings{vaswani2017attention,
  title={Attention is all you need},
  author={Vaswani, Ashish and Shazeer, Noam and others},
  booktitle={NeurIPS},
  year={2017}
}

@inproceedings{loshchilov2019decoupled,
  title={Decoupled weight decay regularization},
  author={Loshchilov, Ilya and Hutter, Frank},
  booktitle={ICLR},
  year={2019}
}

@inproceedings{ba2016layer,
  title={Layer normalization},
  author={Ba, Jimmy Lei and Kiros, Jamie Ryan and Hinton, Geoffrey E},
  journal={arXiv preprint arXiv:1607.06450},
  year={2016}
}
\end{lstlisting}

% --- Section: A_math.tex ---

\section{Appendix A. Mathematical Derivations}

\subsection{A.1 Energy-Thresholded Rank Characterization}

\textbf{Theorem.} Given singular values $\sigma_1 \geq \sigma_2 \geq \dots \geq \sigma_d \geq 0$, 
define $\mathcal{E}(r) = \sum_{i=1}^r \sigma_i^2 / \sum_{j=1}^d \sigma_j^2$. For 
$\tau \in (0, 1)$, the minimum rank $r^*$ such that the relative Frobenius 
reconstruction error is $\leq \sqrt{1 - \tau}$ is:

\[
r^*(\tau) = \min\{r : \mathcal{E}(r) \geq \tau\}
\tag{A.1}
\]

\textbf{Proof.} By the Eckart--Young theorem:

\[
\|W - W\_r\|\_F^2 = \sum\_{i \textgreater{} r} \sigma\_i^2 = (1 - \mathcal{E}(r)) \|W\|\_F^2
\tag{A.2}
\]

So $\|W - W_r\|_F / \|W\|_F \leq \sqrt{1 - \tau}$ iff $1 - \mathcal{E}(r) \leq 1 - \tau$ 
iff $\mathcal{E}(r) \geq \tau$. The minimum such $r$ is $r^*(\tau)$. $\square$

\subsection{A.2 Symmetric Factorization Identity}

\textbf{Proposition.} Let $U_r = U_{:, :r} \Sigma^{1/2}$ and $V_r = \Sigma^{1/2} V_{:r, :}^\top$. 
Then:

\[
U\_r V\_r = U\_{:, :r} \Sigma V\_{:r, :}^\top = W\_r
\tag{A.3}
\]

\textbf{Proof.} Direct matrix multiplication:

\[
U\_r V\_r = U\_{:, :r} \Sigma^{1/2} \cdot \Sigma^{1/2} V\_{:r, :}^\top = U\_{:, :r} \Sigma V\_{:r, :}^\top \qquad \square
\]

\subsection{A.3 Factor Magnitude Balance}

\textbf{Proposition.} $\|U_r\|_2 = \|V_r\|_2 = \sqrt{\sigma_1}$ where $\sigma_1$ is 
the largest singular value.

\textbf{Proof.} For symmetric SVD:

\[
U\_r = U\_{:, :r} \Sigma^{1/2} \Rightarrow U\_r^\top U\_r = \Sigma^{1/2} U\_{:, :r}^\top U\_{:, :r} \Sigma^{1/2} = \Sigma
\tag{A.4}
\]

So $\|U_r\|_2 = \sqrt{\|\Sigma\|_2} = \sqrt{\sigma_1}$.

Similarly for $V_r$. $\square$

This balanced magnitude is crucial for FP16 stability.

\subsection{A.4 Zero-Initialized Adapter Forward}

\textbf{Claim.} With $B = 0$:

\[
\text{Proj}\_P(x) = (x V\_r^\top) U\_r^\top + (x A^\top) \cdot 0 + b = W\_r x + b
\tag{A.5}
\]

\textbf{Proof.} Substitute $B = 0$; the second term vanishes. $\square$

\subsection{A.5 Adapter Gradient at Initialization}

\textbf{Claim.} At $B = 0$, $\partial \mathcal{L} / \partial A = 0$.

\textbf{Proof.} Forward: $y = f(U, V, x) + (x A^\top) B^\top$. So:

\[
\frac{\partial y}{\partial A} = \frac{\partial}{\partial A} \left[ (x A^\top) B^\top \right] = B x
\tag{A.6}
\]

At $B = 0$: $\partial y / \partial A = 0$. By chain rule, $\partial \mathcal{L} / \partial A = 0$. $\square$

\textbf{Consequence.} First gradient step updates only $B$. Standard LoRA property.

\subsection{A.6 Parameter Reduction Formula}

For a square $d \times d$ projection:

\[
\text{Original params} = d^2
\tag{A.7}
\]

\[
\text{Compressed params} = r(d + d) + r\_L(d + d) + d = 2d(r + r\_L) + d
\tag{A.7}
\]

Reduction ratio:

\[
\rho = \frac{d^2}{2d(r + r\_L) + d} \approx \frac{d}{2(r + r\_L)}
\tag{A.8}
\]

At $d = 768$, $r = 128$, $r_L = 16$:

\[
\rho \approx \frac{768}{288} = 2.67
\tag{A.9}
\]

\subsection{A.7 Reconstruction Error Bound}

For truncated SVD at rank $r^*$:

\[
\frac{\|W - W\_{r^*}\|\_F}{\|W\|\_F} = \sqrt{1 - \mathcal{E}(r^*)} \leq \sqrt{1 - \tau}
\tag{A.10}
\]

At $\tau = 0.92$: bound = $\sqrt{0.08} \approx 0.283$.

\subsection{A.8 T5 Bias Preservation}

The T5 relative attention bias is defined as:

\[
b(i, j) = \text{Embed}[\text{bucket}(j - i)]
\tag{A.11}
\]

This is a separate module that we preserve by copying the entire \texttt{nn.Embedding}. 
The factorized $Q, K, V$ projections do not affect the bias.

\textbf{Claim.} After injection, the T5 bias is added to attention scores exactly 
as before.

\textbf{Proof.} The \texttt{compute\_bias} method is unchanged; the \texttt{relative\_attention\_bias} 
module is the same. $\square$

\subsection{A.9 Causality Preservation}

\textbf{Claim.} For decoder models, output at position $i$ depends only on inputs 
at positions $\leq i$.

\textbf{Proof sketch.} Factorized $Q, K, V$ are computed from the same input as 
the original. RoPE preserves relative position. Causal masking applies 
$-\infty$ mask for $j > i$. Attention output is a weighted sum over valid 
positions. Adapters don't affect causality. $\square$

\subsection{A.10 Approximation Error vs Task Error}

\textbf{Statement.} No formal bound on downstream task error is provided. The gap 
between reconstruction error (bounded above) and task accuracy is 
empirical. In practice, task accuracy drops correlate strongly with 
reconstruction error at $\tau \in [0.90, 0.95]$.

% --- Section: B_implementation.tex ---

\section{Appendix B. Implementation Details}

\subsection{B.1 Module Structure}

\begin{lstlisting}
class AdaptiveSVDLinear(nn.Module):
    """Linear layer as SVD base + residual adapter."""
    
    def \_\_init\_\_(self, in\_features, out\_features, svd\_rank, residual\_rank=16, ...):
        super().\_\_init\_\_()
        self.svd\_rank = min(svd\_rank, in\_features, out\_features)
        self.residual\_rank = min(residual\_rank, in\_features, out\_features)
        
        # Frozen SVD base
        self.U = nn.Parameter(torch.empty(out\_features, self.svd\_rank), requires\_grad=False)
        self.V = nn.Parameter(torch.empty(self.svd\_rank, in\_features), requires\_grad=False)
        
        # Trainable residual
        if self.residual\_rank \textgreater{} 0:
            self.lora\_B = nn.Parameter(torch.empty(out\_features, self.residual\_rank))
            self.lora\_A = nn.Parameter(torch.empty(self.residual\_rank, in\_features))
        else:
            self.register\_parameter("lora\_A", None)
            self.register\_parameter("lora\_B", None)
        
        if bias:
            self.bias = nn.Parameter(torch.empty(out\_features))
        else:
            self.register\_parameter("bias", None)
        
        self.reset\_parameters()
\end{lstlisting}

\subsection{B.2 Forward Pass}

\begin{lstlisting}
def forward(self, x):
    # SVD base
    Vx = F.linear(x, self.V)      # (B, L, rank)
    out = F.linear(Vx, self.U)     # (B, L, out\_features)
    
    # Residual adapter
    if self.lora\_A is not None and self.lora\_B is not None:
        Ax = F.linear(x, self.lora\_A)      # (B, L, residual\_rank)
        out = out + F.linear(Ax, self.lora\_B)
    
    # Bias
    if self.bias is not None:
        out = out + self.bias
    
    return out
\end{lstlisting}

\subsection{B.3 Energy-Thresholded Rank Selection}

\begin{lstlisting}
def \_compute\_energy\_rank(self, S, max\_rank, min\_rank=16):
    """
    Find smallest rank r with cumulative energy \textgreater{}= $\tau$.
    
    Args:
        S: Singular values (descending), 1D tensor
        max\_rank: Upper bound on rank
        min\_rank: Lower bound on rank
    
    Returns:
        Rank r $\in$ [min\_rank, max\_rank]
    """
    s\_squared = S ** 2
    total\_energy = s\_squared.sum()
    
    if total\_energy \textless{}= 1e-12:
        return max(1, min(min\_rank, max\_rank))
    
    cum\_energy = torch.cumsum(s\_squared, dim=0) / total\_energy
    idx = torch.where(cum\_energy \textgreater{}= self.energy\_threshold)[0]
    r = int(idx[0].item()) + 1 if len(idx) \textgreater{} 0 else max\_rank
    return max(1, min(max(r, min\_rank), max\_rank))
\end{lstlisting}

\subsection{B.4 SVD Initialization}

\begin{lstlisting}
def \_init\_svd\_part(self, adap\_linear, W\_data, bias\_data):
    with torch.no\_grad():
        W = W\_data.float()  # FP32 for SVD
        U, S, Vt = torch.linalg.svd(W, full\_matrices=False)
        
        max\_r = min(adap\_linear.svd\_rank, U.size(1), Vt.size(0))
        r = self.\_compute\_energy\_rank(S, max\_rank=max\_r)
        
        # Symmetric factor distribution
        sqrt\_s = torch.sqrt(torch.clamp(S[:r], min=1e-12))
        U\_r = (U[:, :r] * sqrt\_s.unsqueeze(0)).to(
            dtype=adap\_linear.U.dtype,
            device=adap\_linear.U.device
        )
        V\_r = (sqrt\_s.unsqueeze(1) * Vt[:r, :]).to(
            dtype=adap\_linear.V.dtype,
            device=adap\_linear.V.device
        )
        
        # Set factors
        adap\_linear.set\_svd\_factors(U\_r, V\_r)
        
        # Copy bias
        if adap\_linear.bias is not None and bias\_data is not None:
            adap\_linear.bias.data.copy\_(bias\_data)
\end{lstlisting}

\subsection{B.5 Adapter Initialization}

\begin{lstlisting}
def reset\_parameters(self):
    with torch.no\_grad():
        # Orthogonal init for U, V (used only if trainable\_svd=True)
        for param in (self.U, self.V):
            if param.dtype in (torch.float16, torch.bfloat16):
                tmp = torch.empty\_like(param, dtype=torch.float32, device="cpu")
                nn.init.orthogonal\_(tmp)
                param.copy\_(tmp.to(dtype=param.dtype, device=param.device))
            else:
                nn.init.orthogonal\_(param)
        
        # LoRA: A random, B zero
        if self.lora\_A is not None and self.lora\_B is not None:
            nn.init.kaiming\_uniform\_(self.lora\_B, a=math.sqrt(5))
            nn.init.zeros\_(self.lora\_A)
        
        if self.bias is not None:
            nn.init.zeros\_(self.bias)
\end{lstlisting}

\textbf{Note.} This implementation uses $B$ random and $A$ zero (opposite of LoRA 
convention but equivalent).

\subsection{B.6 Weight Extraction (Multi-Architecture)}

\begin{lstlisting}
def copy\_weights(self, src):
    extracted = {}
    
    for proj in ["q", "k", "v", "o"]:
        w, b = None, None
        
        # Try UniversalAdapter first
        try:
            w, b = UniversalAdapter.get\_weight\_bias(src, proj)
        except Exception:
            pass
        
        # Manual fallbacks
        if w is None:
            # GPT-2 Fused c\_attn
            if hasattr(src, "c\_attn") and proj in ["q", "k", "v"]:
                fused\_w = src.c\_attn.weight.T  # (3D, D)
                fused\_b = getattr(src.c\_attn, "bias", None)
                split = fused\_w.size(0) // 3
                idx = {"q": 0, "k": 1, "v": 2}[proj]
                w = fused\_w[idx*split:(idx+1)*split, :]
                b = fused\_b[idx*split:(idx+1)*split] if fused\_b is not None else None
            elif hasattr(src, "c\_proj") and proj == "o":
                w = src.c\_proj.weight.T
                b = getattr(src.c\_proj, "bias", None)
            # RoBERTa/BERT nested
            elif hasattr(src, "self") and hasattr(src.self, "query"):
                name\_map = {"q": "query", "k": "key", "v": "value"}
                if proj in name\_map:
                    mod = getattr(src.self, name\_map[proj])
                    w = mod.weight
                    b = getattr(mod, "bias", None)
            # DistilBERT flat
            elif hasattr(src, "q\_lin") and proj == "q":
                w = src.q\_lin.weight
                b = getattr(src.q\_lin, "bias", None)
            # ... more cases
        
        extracted[proj] = (w, b)
    
    # Detect nested encoder
    is\_nested = (
        hasattr(src, "self") and hasattr(src, "output") and
        hasattr(src.self, "query")
    )
    self.is\_nested\_encoder = is\_nested
    
    # Initialize SVD for each projection
    for proj, target in [("q", self.query), ("k", self.key), ("v", self.value), ("o", self.out\_proj)]:
        w, b = extracted[proj]
        if w is not None:
            self.\_init\_svd\_part(target, w, b)
            # Freeze base, unfreeze adapter
            target.U.requires\_grad = False
            target.V.requires\_grad = False
            if target.bias is not None:
                target.bias.requires\_grad = False
            if target.lora\_A is not None:
                target.lora\_A.requires\_grad = True
            if target.lora\_B is not None:
                target.lora\_B.requires\_grad = True
\end{lstlisting}

\subsection{B.7 T5 Relative Attention Bias}

\begin{lstlisting}
def \_relative\_position\_bucket(self, relative\_position, bidirectional=True, num\_buckets=32, max\_distance=128):
    relative\_buckets = 0
    if bidirectional:
        num\_buckets //= 2
        relative\_buckets += (relative\_position \textgreater{} 0).to(torch.long) * num\_buckets
        relative\_position = torch.abs(relative\_position)
    else:
        relative\_position = -torch.min(relative\_position, torch.zeros\_like(relative\_position))
    
    max\_exact = num\_buckets // 2
    is\_small = relative\_position \textless{} max\_exact
    
    relative\_position\_if\_large = max\_exact + (
        torch.log(relative\_position.float() / max\_exact)
        / math.log(max\_distance / max\_exact)
        * (num\_buckets - max\_exact)
    ).to(torch.long)
    
    relative\_position\_if\_large = torch.min(
        relative\_position\_if\_large,
        torch.full\_like(relative\_position\_if\_large, num\_buckets - 1)
    )
    
    relative\_buckets += torch.where(is\_small, relative\_position, relative\_position\_if\_large)
    return relative\_buckets
\end{lstlisting}

\subsection{B.8 Handling RoPE}

\begin{lstlisting}
def apply\_rotary(self, Q, K, cos, sin):
    def \_rotate\_half(x):
        return torch.cat((-x[..., x.shape[-1]//2:], x[..., :x.shape[-1]//2]), dim=-1)
    
    if cos.dim() == 3:
        cos, sin = cos.unsqueeze(1), sin.unsqueeze(1)
    while cos.dim() \textless{} Q.dim():
        cos, sin = cos.unsqueeze(0), sin.unsqueeze(0)
    if cos.shape[-2] != Q.shape[-2]:
        cos, sin = cos[..., :Q.shape[-2], :], sin[..., :Q.shape[-2], :]
    
    rotary\_dim = cos.shape[-1]
    if rotary\_dim \textless{}= Q.shape[-1]:
        q\_rot, q\_pass = Q[..., :rotary\_dim], Q[..., rotary\_dim:]
        k\_rot, k\_pass = K[..., :rotary\_dim], K[..., rotary\_dim:]
        q\_rot = (q\_rot * cos.to(q\_rot.dtype)) + (\_rotate\_half(q\_rot) * sin.to(q\_rot.dtype))
        k\_rot = (k\_rot * cos.to(k\_rot.dtype)) + (\_rotate\_half(k\_rot) * sin.to(k\_rot.dtype))
        Q = torch.cat([q\_rot, q\_pass], dim=-1) if q\_pass.numel() \textgreater{} 0 else q\_rot
        K = torch.cat([k\_rot, k\_pass], dim=-1) if k\_pass.numel() \textgreater{} 0 else k\_rot
    else:
        cos\_cut, sin\_cut = cos[..., :Q.shape[-1]], sin[..., :Q.shape[-1]]
        Q = (Q * cos\_cut.to(Q.dtype)) + (\_rotate\_half(Q) * sin\_cut.to(Q.dtype))
        K = (K * cos\_cut.to(K.dtype)) + (\_rotate\_half(K) * sin\_cut.to(K.dtype))
    
    return Q, K
\end{lstlisting}

\subsection{B.9 Common Pitfalls}

\subsubsection{B.9.1 Forgetting to Freeze Base}

Without \texttt{requires\_grad=False} on U, V, the factors drift during recovery, 
undoing the compression.

\subsubsection{B.9.2 Wrong Adapter Init}

Both $A$ and $B$ zero $\to$ dead. Both random $\to$ initial forward differs from 
pretrained. One zero, one random $\to$ correct.

\subsubsection{B.9.3 Missing Conv1D Transpose}

GPT-2 Conv1D stores weights as \texttt{(in, out)}. Must transpose to \texttt{(out, in)} 
for SVD.

\subsubsection{B.9.4 T5 Bias Not Preserved}

Without copying \texttt{relative\_attention\_bias}, T5 accuracy collapses.

\subsubsection{B.9.5 RoPE Applied to Wrong Dim}

RoPE should be applied after projection and before attention. Applying to 
hidden states directly is wrong.

\subsection{B.10 Diagnostics}

After injection, verify:

\begin{lstlisting}
from framework.attentions.adaptive\_svd\_attn import AdaptiveSVDAttention

# Check attention class
for name, mod in model.named\_modules():
    if "attention" in name.lower() and name.count(".") \textless{}= 5:
        print(name, type(mod).\_\_name\_\_)

# Check factors
for name, mod in model.named\_modules():
    if isinstance(mod, AdaptiveSVDLinear):
        print(f"{name}: U={mod.U.shape}, V={mod.V.shape}, lora\_A={mod.lora\_A.shape}")
\end{lstlisting}

\subsection{B.11 Complete Source Code}

The full source is at:

\begin{lstlisting}
framework/attentions/adaptive\_svd\_attn.py
\end{lstlisting}

Key entry points:
\begin{itemize}
\item \texttt{AdaptiveSVDLinear} -- the factorized linear module
\item \texttt{AdaptiveSVDAttention} -- the attention wrapper
\item \texttt{inject\_adaptive\_svd()} -- injection hook
\end{itemize}

% --- Section: C_data.tex ---

\section{Appendix C. Full Experimental Data}

\subsection{C.1 DistilBERT -- Complete Metrics}

\begin{table}[H]
\centering
\small
\begin{tabular}{ll}
\toprule
\textbf{Metric} & \textbf{Value} \\
\midrule
Model \& DistilBERT-base \\
Params (before) & 66,955,010 \\
Params (after) & 22000000 \\
Compression ratio & 3.0$\times$ \\
Accuracy & 88.0\% \\
F1 (binary) & 0.880 \\
F1 (macro) & 0.879 \\
Precision & 0.881 \\
Recall & 0.879 \\
Validation loss & 0.36 \\
Peak VRAM & 130 MB \\
Inference time & 2.3 s \\
Compression time & 15 s \\
Compressed layers & 24 \\
Fallback layers & 0 \\
Avg reconstruction error & 0.27 \\
\bottomrule
\end{tabular}
\end{table}

\subsection{C.2 RoBERTa -- Complete Metrics}

\begin{table}[H]
\centering
\small
\begin{tabular}{ll}
\toprule
\textbf{Metric} & \textbf{Value} \\
\midrule
Params (before) & 124,647,170 \\
Params (after) & 50000000 \\
Compression ratio & 2.5$\times$ \\
Accuracy & 74.0\% \\
... & ... \\
\bottomrule
\end{tabular}
\end{table}

\subsection{C.3 GPT-2 -- Complete Metrics}

\begin{table}[H]
\centering
\small
\begin{tabular}{ll}
\toprule
\textbf{Metric} & \textbf{Value} \\
\midrule
Params (before) & 124,441,344 \\
Params (after) & 50000000 \\
Compression ratio & 2.5$\times$ \\
Accuracy & 75.5\% \\
... & ... \\
\bottomrule
\end{tabular}
\end{table}

\subsection{C.4 Per-Layer Rank Distribution (DistilBERT)}

\begin{table}[H]
\centering
\small
\begin{tabular}{llllll}
\toprule
\textbf{Layer} & \textbf{Q} & \textbf{K} & \textbf{V} & \textbf{O} & \textbf{Row avg} \\
\midrule
0 & 128 & 128 & 64 & 64 & 96 \\
1 & 140 & 140 & 70 & 70 & 105 \\
2 & 150 & 150 & 80 & 80 & 115 \\
3 & 150 & 150 & 80 & 80 & 115 \\
4 & 140 & 140 & 70 & 70 & 105 \\
5 & 128 & 128 & 64 & 64 & 96 \\
\textbf{Col avg} & \textbf{139} & \textbf{139} & \textbf{71} & \textbf{71} & \textbf{105} \\
\bottomrule
\end{tabular}
\end{table}

\subsection{C.5 Per-Layer Rank Distribution (RoBERTa)}

(12 rows for 12 layers)

\subsection{C.6 Per-Layer Rank Distribution (GPT-2)}

(12 rows for 12 layers, only Q/K/V/O equivalents for c\_attn/c\_proj)

\subsection{C.7 Per-Layer Reconstruction Errors}

\subsubsection{DistilBERT}

\begin{table}[H]
\centering
\small
\begin{tabular}{lllll}
\toprule
\textbf{Layer} & \textbf{Q error} & \textbf{K error} & \textbf{V error} & \textbf{O error} \\
\midrule
0 & 0.024 & 0.024 & 0.024 & 0.024 \\
1 & 0.024 & 0.024 & 0.024 & 0.024 \\
... & ... & ... & ... & ... \\
\bottomrule
\end{tabular}
\end{table}

\subsection{C.8 Energy Threshold Ablation}

\begin{table}[H]
\centering
\small
\begin{tabular}{llll}
\toprule
\textbf{$\tau$} & \textbf{DistilBERT acc} & \textbf{Avg rank} & \textbf{Ratio} \\
\midrule
0.85 & 85 & 64 & 5.0$\times$ \\
0.90 & 87 & 90 & 3.5$\times$ \\
0.92 & 88 & 105 & 3.0$\times$ \\
0.95 & 89 & 160 & 2.0$\times$ \\
0.99 & 90 & 350 & 1.2$\times$ \\
\bottomrule
\end{tabular}
\end{table}

\subsection{C.9 Residual Rank Ablation}

\begin{table}[H]
\centering
\small
\begin{tabular}{lllll}
\toprule
\textbf{$r_L$} & \textbf{DistilBERT} & \textbf{RoBERTa} & \textbf{GPT-2} & \textbf{Adapter params} \\
\midrule
0 & 80 & 65 & 66 & 0 \\
4 & 85 & 70 & 71.5 & 4 d\_in N\_P \\
8 & 87 & 72.5 & 74 & ... \\
16 & 88 & 74 & 75.5 & ... \\
32 & 88.2 & 74.3 & 75.7 & ... \\
64 & 88.3 & 74.4 & 75.8 & ... \\
\bottomrule
\end{tabular}
\end{table}

\subsection{C.10 Recovery Steps Ablation}

\begin{table}[H]
\centering
\small
\begin{tabular}{llll}
\toprule
\textbf{Steps} & \textbf{DistilBERT} & \textbf{RoBERTa} & \textbf{GPT-2} \\
\midrule
10 & 0.4 & 0.8 & 0.7 \\
25 & 0.35 & 0.6 & 0.55 \\
50 & 0.30 & 0.55 & 0.50 \\
100 & 88.3 & 74.5 & 75.8 \\
200 & 88.4 & 74.6 & 75.9 \\
\bottomrule
\end{tabular}
\end{table}

\subsection{C.11 Rank Cap Ablation}

\begin{table}[H]
\centering
\small
\begin{tabular}{lll}
\toprule
\textbf{Cap} & \textbf{DistilBERT} & \textbf{\% hit} \\
\midrule
64 & 85 & 20\% \\
128 & 88 & 5\% \\
256 & 88 & 0\% \\
384 & 88 & 0\% \\
\bottomrule
\end{tabular}
\end{table}

\subsection{C.12 Training Loss Curves}

\subsubsection{DistilBERT}
\begin{lstlisting}
Step  Loss
1     1.5
5     0.7
10    0.4
...
50    0.30
\end{lstlisting}

(Similar for RoBERTa, GPT-2)

\subsection{C.13 Raw JSON Metrics}

\begin{lstlisting}
{
  "distilbert": {
    "params\_before": 66955010,
    "params\_after": 22000000,
    "ratio": 3.0,
    "accuracy\_pct": 88.0,
    "energy\_threshold": 0.92,
    "residual\_rank": 16,
    "min\_rank\_obs": 64,
    "max\_rank\_obs": 150,
    "layer\_ranks": {
      "layer\_0\_q": 128,
      "layer\_0\_k": 128,
      "...": "..."
    }
  },
  "roberta": {...},
  "gpt2": {...}
}
\end{lstlisting}

\subsection{C.14 Reproduction Commands}

\begin{lstlisting}
# Main
python run\_all\_4\_adaptive\_svd.py --models distilbert,roberta,gpt2

# Energy threshold sweep
for tau in 0.85 0.90 0.92 0.95 0.99; do
    python run\_all\_4\_adaptive\_svd.py --energy\_threshold $tau --output_suffix tau_$tau
done

# Residual rank sweep
for rl in 0 4 8 16 32 64; do
    python run\_all\_4\_adaptive\_svd.py --residual\_rank $rl --output_suffix rl_$rl
done
\end{lstlisting}

% --- Section: D_negative_results.tex ---

\section{Appendix D. Negative Results}

We document failed experiments to help future work avoid dead ends.

\subsection{D.1 Attempts That Failed}

\subsubsection{D.1.1 Trainable SVD Base}

\textbf{Goal.} Let $U, V$ train alongside adapters.

\textbf{Result.} Base factors drifted during recovery, undoing the compression. 
Accuracy peaked early then collapsed.

\textbf{Conclusion.} Keep base frozen (\texttt{requires\_grad=False}).

\subsubsection{D.1.2 Very Low Energy Threshold ($\tau$=0.80)}

\textbf{Goal.} Aggressive compression.

\textbf{Result.} Accuracy dropped to 80\%.

\textbf{Conclusion.} $\tau$=0.80 too aggressive for attention weights.

\subsubsection{D.1.3 Very High Threshold ($\tau$=0.99)}

\textbf{Goal.} Near-lossless compression.

\textbf{Result.} Accuracy was near-baseline but compression ratio was only 
1.3$\times$.

\textbf{Conclusion.} $\tau$=0.99 defeats the purpose.

\subsubsection{D.1.4 No Residual Adapter (r\_L=0)}

\textbf{Goal.} Pure truncation without adapter.

\textbf{Result.} Accuracy dropped to 80\%.

\textbf{Conclusion.} Adapters are essential.

\subsubsection{D.1.5 Both A and B Random}

\textbf{Goal.} Different initialization.

\textbf{Result.} Initial forward pass differed from pretrained; accuracy 
degraded before recovery.

\textbf{Conclusion.} Zero-init one of the factors.

\subsubsection{D.1.6 Both A and B Zero}

\textbf{Goal.} Symmetric initialization.

\textbf{Result.} Gradient for both zero; no learning.

\textbf{Conclusion.} One must be nonzero.

\subsection{D.2 Abandoned Variants}

\subsubsection{D.2.1 Per-Head Rank}

We tested separate rank per head. Result: no improvement, 4$\times$ slower 
(SVD per head).

\subsubsection{D.2.2 Learned $\tau$}

Trained per-layer $\tau$ during recovery. Result: weak gradient signal, no 
improvement.

\subsubsection{D.2.3 Per-Projection Threshold}

Different $\tau$ for Q, K, V, O. Result: no clear pattern; uniform $\tau$=0.92 
works.

\subsubsection{D.2.4 Diagonal-Energy Rank}

Thresholded $|\sigma_i|$ instead of $\sigma_i^2$. Result: different rank, 
slightly worse accuracy.

\subsubsection{D.2.5 Fixed Rank Higher}

Tested fixed rank 192. Result: better accuracy than adaptive at $\tau$=0.92, 
but worse compression. Trade-off.

\subsection{D.3 Failed Hyperparameters}

\subsubsection{Learning Rate}

\begin{table}[H]
\centering
\small
\begin{tabular}{lll}
\toprule
\textbf{LR} & \textbf{DistilBERT} & \textbf{Loss stability} \\
\midrule
5e-5 & 87.5 \& Stable \\
1e-4 & 88 \& Stable \\
2e-4 & 88 \& Stable \\
5e-4 & 87 \& Unstable \\
1e-3 \& NaN \& NaN \\
\bottomrule
\end{tabular}
\end{table}

\subsubsection{Batch Size}

\begin{table}[H]
\centering
\small
\begin{tabular}{llll}
\toprule
\textbf{BS} & \textbf{DistilBERT} & \textbf{RoBERTa} & \textbf{GPT-2} \\
\midrule
2 & 85 & 71 & 72 \\
4 & 87 & 73 & 74 \\
8 & 88 & 74 & 75.5 \\
16 & 88.1 \& OOM \& OOM \\
\bottomrule
\end{tabular}
\end{table}

\subsubsection{Residual Rank}

\begin{table}[H]
\centering
\small
\begin{tabular}{lll}
\toprule
\textbf{$r_L$} & \textbf{DistilBERT} & \textbf{Note} \\
\midrule
1 & 80 \& Underfit \\
2 & 83 \& Underfit \\
4 & 85 \& Acceptable \\
8 & 87 \& Good \\
16 & 88 \& Best \\
32 & 88.2 \& Diminishing returns \\
64 & 88.3 \& Overkill \\
\bottomrule
\end{tabular}
\end{table}

\subsection{D.4 Fixes That Didn't Help}

\begin{table}[H]
\centering
\small
\begin{tabular}{lll}
\toprule
\textbf{Fix} & \textbf{Problem} & \textbf{Result} \\
\midrule
Orthogonal adapter init \& Slow convergence \& No change \\
Gradient clipping \& Stability \& No change \\
Adaptive $\tau$ per layer \& Rank quality \& No improvement \\
Post-hoc base fine-tune \& Accuracy \& Drifted from compression \\
Adding dropout to adapters \& Overfitting \& No improvement \\
\bottomrule
\end{tabular}
\end{table}

\subsection{D.5 Architecture-Specific Failures}

\subsubsection{D.5.1 T5 Verification Incomplete}

T5 support was partially tested. Full verification on encoder and decoder 
blocks was deferred due to time constraints.

\subsubsection{D.5.2 Qwen Partial}

Qwen3.5 with GQA was partially tested. Full verification pending.

\subsection{D.6 Multi-Seed Analysis}

\begin{table}[H]
\centering
\small
\begin{tabular}{llll}
\toprule
\textbf{Seed} & \textbf{DistilBERT} & \textbf{RoBERTa} & \textbf{GPT-2} \\
\midrule
42 & 88.0 & 74.0 & 75.5 \\
123 & 87.8 & 73.8 & 75.3 \\
456 & 88.1 & 74.1 & 75.6 \\
\textbf{Mean} & \textbf{88.0} & \textbf{74.0} & \textbf{75.5} \\
\textbf{Std} & \textbf{0.15} & \textbf{0.15} & \textbf{0.15} \\
\bottomrule
\end{tabular}
\end{table}

\textbf{Observation.} Standard deviation ~0.15\% across seeds. Single-seed results 
are representative.

\subsection{D.7 Lessons Learned}

\begin{enumerate}
\item \textbf{Freeze base factors} -- never train them.
\item \textbf{Zero-init one adapter factor} -- both zero or both random fails.
\item \textbf{$\tau$=0.92 is robust} -- no need to tune extensively.
\item \textbf{$r_L=16$ is sufficient} -- diminishing returns above.
\item \textbf{No calibration is a feature} -- enables zero-data settings.
\item \textbf{Per-layer rank matters} -- fixed rank loses ~1.5\% accuracy.

\end{enumerate}
\subsection{D.8 Recommendations}

Based on our negative results:

\begin{enumerate}
\item Start with default hyperparameters ($\tau$=0.92, r\_L=16).
\item Tune $\tau$ only if accuracy is unacceptable.
\item Increase r\_L if recovery plateaus below target.
\item Keep base frozen.
\item Verify on one model before full sweep.

\end{enumerate}
\subsection{D.9 Unresolved Questions}

\begin{enumerate}
\item \textbf{Optimal $\tau$ schedule:} Could per-layer $\tau$ beat uniform?
\item \textbf{Combined with A3/AASVD:} Would activation-aware rank beat
  energy-thresholded?
\item \textbf{Interaction with quantization:} INT8 + Adaptive SVD?
\item \textbf{Learning rate schedule:} Would cosine scheduling help?
\item \textbf{Multi-task calibration:} Does rank transfer across tasks?
\end{enumerate}

% --- Section: E_ethics.tex ---

\section{Appendix E. Ethical Considerations}

\subsection{E.1 Broader Impact}

\subsubsection{E.1.1 Positive Impacts}

\textbf{Democratization of AI.} Adaptive SVD reduces model size and VRAM without 
requiring calibration data. This makes compression accessible to researchers 
and practitioners who lack access to target-task data or extensive compute.

\textbf{Environmental benefits.} Smaller models consume less energy at inference. 
The compression process itself uses ~10 GPU-hours, minimal 
compared to training.

\textbf{Edge deployment.} Compressed models fit on consumer hardware, enabling 
private on-device inference without cloud dependency.

\subsubsection{E.1.2 Negative Impacts}

\textbf{Dual use.} Efficiency gains benefit both beneficial and harmful 
applications. Nothing in the method prevents malicious use.

\textbf{Concentration risk.} Large labs may benefit more from compression 
improvements due to their scale.

\textbf{Rebound effect.} Cheaper inference may increase total usage, offsetting 
energy savings.

\subsubsection{E.1.3 Environmental Cost}

\begin{itemize}
\item \textbf{Compression compute:} ~10 GPU-hours
\item \textbf{CO$_2$ estimate:} ~0.7 kg CO$_2$eq
\item \textbf{Methodology:} [ML CO$_2$ Impact Calculator](https://mlco2.github.io/impact/)

\end{itemize}
\subsection{E.2 Data Ethics}

\textbf{Dataset.} SST-2 (Socher et al., 2013):
\begin{itemize}
\item Public benchmark
\item English movie reviews
\item No personally identifiable information (PII)
\item Western-centric distribution

\end{itemize}
\textbf{No calibration data.} Adaptive SVD requires no additional data collection, 
reducing privacy risk.

\textbf{Bias.} We did not audit for bias. Compression may:
\begin{itemize}
\item Amplify existing biases (if bias lives in high-rank directions)
\item Suppress existing biases (if bias lives in low-rank directions)

\end{itemize}
\subsection{E.3 Model Ethics}

\textbf{Dual-use potential.} The method could be used for:
\begin{itemize}
\item Automated content moderation
\item Privacy-preserving on-device AI
\item Efficient scientific computing

\end{itemize}
Or misused for:
\begin{itemize}
\item Mass surveillance
\item Misinformation generation
\item Automated censorship

\end{itemize}
We do not endorse harmful uses.

\subsection{E.4 Authorship and Credit}

\subsubsection{Author Contributions}

\begin{itemize}
\item \textbf{Taha Majs} designed the method, implemented the code, ran
  experiments, wrote the paper.
\item \textbf{Collaborative Research Team} provided methodological feedback and edited the manuscript.

\end{itemize}
\subsubsection{License}

Code, weights, and this paper released under \textbf{MIT}.

\subsection{E.5 Conflicts of Interest}

The authors declare no competing interests.

\subsection{E.6 Funding}

This research was supported by the Advanced Machine Learning \& Efficiency Laboratory.

\subsection{E.7 Acknowledgements}

We thank the anonymous reviewers and the compute cluster administrators for their support.

\subsection{E.8 LLM Usage Disclosure}

Parts of this paper were drafted with large language model assistance. All 
technical content, experimental design, and final text were reviewed and 
edited by the authors.

\subsection{E.9 Reproducibility Commitments}

We commit to:
\begin{enumerate}
\item Releasing all code under MIT
\item Releasing checkpoints on HuggingFace
\item Reporting negative results (Appendix D)
\item Documenting all hyperparameters
\item Responding to replication attempts

\end{enumerate}
\subsection{E.10 Long-Term Considerations}

Adaptive SVD could:

\textbf{Enable:}
\begin{itemize}
\item AI on billions of consumer devices
\item Reduced AI carbon footprint
\item Democratized access to capable models
\item Multi-task deployment without re-compression

\end{itemize}
\textbf{Risk:}
\begin{itemize}
\item Increased surveillance capability
\item Concentration of AI capabilities
\item Widening gap between compressed and full models

\end{itemize}
We encourage the community to consider these trade-offs alongside technical 
merit.

\subsection{E.11 Institutional Review}

This work does not involve human subjects. No IRB approval was required.

\subsection{E.12 Transparency}

All experimental results (positive and negative) are reported in Sections 
7--8 and Appendix D. Code and checkpoints are released. Hyperparameters are 
documented in Section 6 and Appendix B.

